% Môn: Vật lý
% Lớp: 10
% Chương: Chương III. Động lực học
% Bài: L10C3 Bài 16. Định luật 3 Newton
% Số câu: 162
% App đọc file này trực tiếp — sửa rồi Commit trên GitHub, bấm Đồng bộ GitHub .tex

% L10C3 Bài 16. Định luật 3 Newton

% ===== Câu 1 =====
% ID: AIQ_L10C3_FULL_0487
% Mức: NB
\dangbt{Nhận biết lực và phản lực theo định luật III Newton}
\begin{ex}
% ID: AIQ_L10C3_FULL_0487
% Mức: NB
Vật $A$ tác dụng lên vật $B$ lực $5\,\text{N}$. Theo định luật III Newton, độ lớn lực $B$ tác dụng lại $A$ bằng bao nhiêu?
\choice
{\True 5 $N$}
{6 $N$}
{4 $N$}
{10 $N$}
\loigiai{
Chọn A. Hai lực tương tác có cùng độ lớn nên lực $B$ tác dụng lên $A$ bằng $5\,\text{N}$.
}
\end{ex}

% ===== Câu 2 =====
% ID: AIQ_L10C3_FULL_0488
% Mức: NB
\begin{ex}
% ID: AIQ_L10C3_FULL_0488
% Mức: NB
Vật $A$ tác dụng lên vật $B$ lực $7\,\text{N}$. Theo định luật III Newton, độ lớn lực $B$ tác dụng lại $A$ bằng bao nhiêu?
\choice
{7 $N$}
{\True 7 $N$}
{6 $N$}
{14 $N$}
\loigiai{
Chọn B. Hai lực tương tác có cùng độ lớn nên lực $B$ tác dụng lên $A$ bằng $7\,\text{N}$.
}
\end{ex}

% ===== Câu 3 =====
% ID: AIQ_L10C3_FULL_0489
% Mức: NB
\begin{ex}
% ID: AIQ_L10C3_FULL_0489
% Mức: NB
Vật $A$ tác dụng lên vật $B$ lực $9\,\text{N}$. Theo định luật III Newton, độ lớn lực $B$ tác dụng lại $A$ bằng bao nhiêu?
\choice
{9 $N$}
{10 $N$}
{\True 9 $N$}
{18 $N$}
\loigiai{
Chọn C. Hai lực tương tác có cùng độ lớn nên lực $B$ tác dụng lên $A$ bằng $9\,\text{N}$.
}
\end{ex}

% ===== Câu 4 =====
% ID: AIQ_L10C3_FULL_0490
% Mức: NB
\begin{ex}
% ID: AIQ_L10C3_FULL_0490
% Mức: NB
Vật $A$ tác dụng lên vật $B$ lực $11\,\text{N}$. Theo định luật III Newton, độ lớn lực $B$ tác dụng lại $A$ bằng bao nhiêu?
\choice
{11 $N$}
{12 $N$}
{10 $N$}
{\True 11 $N$}
\loigiai{
Chọn D. Hai lực tương tác có cùng độ lớn nên lực $B$ tác dụng lên $A$ bằng $11\,\text{N}$.
}
\end{ex}

% ===== Câu 5 =====
% ID: AIQ_L10C3_FULL_0491
% Mức: TH
\begin{ex}
% ID: AIQ_L10C3_FULL_0491
% Mức: TH
Vật $A$ tác dụng lên vật $B$ lực $13\,\text{N}$. Theo định luật III Newton, độ lớn lực $B$ tác dụng lại $A$ bằng bao nhiêu?
\choice
{\True 13 $N$}
{14 $N$}
{12 $N$}
{26 $N$}
\loigiai{
Chọn A. Hai lực tương tác có cùng độ lớn nên lực $B$ tác dụng lên $A$ bằng $13\,\text{N}$.
}
\end{ex}

% ===== Câu 6 =====
% ID: AIQ_L10C3_FULL_0492
% Mức: TH
\begin{ex}
% ID: AIQ_L10C3_FULL_0492
% Mức: TH
Vật $A$ tác dụng lên vật $B$ lực $15\,\text{N}$. Theo định luật III Newton, độ lớn lực $B$ tác dụng lại $A$ bằng bao nhiêu?
\choice
{15 $N$}
{\True 15 $N$}
{14 $N$}
{30 $N$}
\loigiai{
Chọn B. Hai lực tương tác có cùng độ lớn nên lực $B$ tác dụng lên $A$ bằng $15\,\text{N}$.
}
\end{ex}

% ===== Câu 7 =====
% ID: AIQ_L10C3_FULL_0493
% Mức: TH
\begin{ex}
% ID: AIQ_L10C3_FULL_0493
% Mức: TH
Vật $A$ tác dụng lên vật $B$ lực $17\,\text{N}$. Theo định luật III Newton, độ lớn lực $B$ tác dụng lại $A$ bằng bao nhiêu?
\choice
{17 $N$}
{18 $N$}
{\True 17 $N$}
{34 $N$}
\loigiai{
Chọn C. Hai lực tương tác có cùng độ lớn nên lực $B$ tác dụng lên $A$ bằng $17\,\text{N}$.
}
\end{ex}

% ===== Câu 8 =====
% ID: AIQ_L10C3_FULL_0494
% Mức: TH
\begin{ex}
% ID: AIQ_L10C3_FULL_0494
% Mức: TH
Vật $A$ tác dụng lên vật $B$ lực $19\,\text{N}$. Theo định luật III Newton, độ lớn lực $B$ tác dụng lại $A$ bằng bao nhiêu?
\choice
{19 $N$}
{20 $N$}
{18 $N$}
{\True 19 $N$}
\loigiai{
Chọn D. Hai lực tương tác có cùng độ lớn nên lực $B$ tác dụng lên $A$ bằng $19\,\text{N}$.
}
\end{ex}

% ===== Câu 9 =====
% ID: AIQ_L10C3_FULL_0495
% Mức: VD
\begin{ex}
% ID: AIQ_L10C3_FULL_0495
% Mức: VD
Vật $A$ tác dụng lên vật $B$ lực $5\,\text{N}$. Theo định luật III Newton, độ lớn lực $B$ tác dụng lại $A$ bằng bao nhiêu?
\choice
{\True 5 $N$}
{6 $N$}
{4 $N$}
{10 $N$}
\loigiai{
Chọn A. Hai lực tương tác có cùng độ lớn nên lực $B$ tác dụng lên $A$ bằng $5\,\text{N}$.
}
\end{ex}

% ===== Câu 10 =====
% ID: AIQ_L10C3_FULL_0496
% Mức: VD
\begin{ex}
% ID: AIQ_L10C3_FULL_0496
% Mức: VD
Vật $A$ tác dụng lên vật $B$ lực $7\,\text{N}$. Theo định luật III Newton, độ lớn lực $B$ tác dụng lại $A$ bằng bao nhiêu?
\choice
{7 $N$}
{\True 7 $N$}
{6 $N$}
{14 $N$}
\loigiai{
Chọn B. Hai lực tương tác có cùng độ lớn nên lực $B$ tác dụng lên $A$ bằng $7\,\text{N}$.
}
\end{ex}

% ===== Câu 11 =====
% ID: AIQ_L10C3_FULL_0497
% Mức: TH
\begin{ex}
% ID: AIQ_L10C3_FULL_0497
% Mức: TH
Vật $A$ tác dụng lên vật $B$ lực $15\,\text{N}$. Theo định luật III Newton, độ lớn lực $B$ tác dụng lại $A$ bằng bao nhiêu?
\choiceTF
{\True Giá trị cần tìm là $15\,\text{N}$.}
{Giá trị cần tìm là $16\,\text{N}$.}
{\True Phải xác định đúng các lực tác dụng và chọn chiều dương trước khi lập phương trình.}
{Có thể bỏ qua điều kiện áp dụng và chiều của lực mà kết quả vẫn luôn đúng.}
\loigiai{
\begin{itemchoice}
\item (Đúng) Giá trị cần tìm là $15\,\text{N}$. (Vì): Hai lực tương tác có cùng độ lớn nên lực $B$ tác dụng lên $A$ bằng $15\,\text{N}$. b) Sai vì sai giá trị. c) Đúng. d) Sai vì phải kiểm tra điều kiện áp dụng và chiều lực
\item (Sai) Giá trị cần tìm là $16\,\text{N}$.
\item (Đúng) Phải xác định đúng các lực tác dụng và chọn chiều dương trước khi lập phương trình.
\item (Sai) Có thể bỏ qua điều kiện áp dụng và chiều của lực mà kết quả vẫn luôn đúng.
\end{itemchoice}
}
\end{ex}

% ===== Câu 12 =====
% ID: AIQ_L10C3_FULL_0498
% Mức: TH
\begin{ex}
% ID: AIQ_L10C3_FULL_0498
% Mức: TH
Vật $A$ tác dụng lên vật $B$ lực $17\,\text{N}$. Theo định luật III Newton, độ lớn lực $B$ tác dụng lại $A$ bằng bao nhiêu?
\choiceTF
{\True Giá trị cần tìm là $17\,\text{N}$.}
{Giá trị cần tìm là $18\,\text{N}$.}
{\True Phải xác định đúng các lực tác dụng và chọn chiều dương trước khi lập phương trình.}
{Có thể bỏ qua điều kiện áp dụng và chiều của lực mà kết quả vẫn luôn đúng.}
\loigiai{
\begin{itemchoice}
\item (Đúng) Giá trị cần tìm là $17\,\text{N}$. (Vì): Hai lực tương tác có cùng độ lớn nên lực $B$ tác dụng lên $A$ bằng $17\,\text{N}$. b) Sai vì sai giá trị. c) Đúng. d) Sai vì phải kiểm tra điều kiện áp dụng và chiều lực
\item (Sai) Giá trị cần tìm là $18\,\text{N}$.
\item (Đúng) Phải xác định đúng các lực tác dụng và chọn chiều dương trước khi lập phương trình.
\item (Sai) Có thể bỏ qua điều kiện áp dụng và chiều của lực mà kết quả vẫn luôn đúng.
\end{itemchoice}
}
\end{ex}

% ===== Câu 13 =====
% ID: AIQ_L10C3_FULL_0499
% Mức: VD
\begin{ex}
% ID: AIQ_L10C3_FULL_0499
% Mức: VD
Vật $A$ tác dụng lên vật $B$ lực $19\,\text{N}$. Theo định luật III Newton, độ lớn lực $B$ tác dụng lại $A$ bằng bao nhiêu?
\choiceTF
{\True Giá trị cần tìm là $19\,\text{N}$.}
{Giá trị cần tìm là $20\,\text{N}$.}
{\True Phải xác định đúng các lực tác dụng và chọn chiều dương trước khi lập phương trình.}
{Có thể bỏ qua điều kiện áp dụng và chiều của lực mà kết quả vẫn luôn đúng.}
\loigiai{
\begin{itemchoice}
\item (Đúng) Giá trị cần tìm là $19\,\text{N}$. (Vì): Hai lực tương tác có cùng độ lớn nên lực $B$ tác dụng lên $A$ bằng $19\,\text{N}$. b) Sai vì sai giá trị. c) Đúng. d) Sai vì phải kiểm tra điều kiện áp dụng và chiều lực
\item (Sai) Giá trị cần tìm là $20\,\text{N}$.
\item (Đúng) Phải xác định đúng các lực tác dụng và chọn chiều dương trước khi lập phương trình.
\item (Sai) Có thể bỏ qua điều kiện áp dụng và chiều của lực mà kết quả vẫn luôn đúng.
\end{itemchoice}
}
\end{ex}

% ===== Câu 14 =====
% ID: AIQ_L10C3_FULL_0500
% Mức: VD
\begin{ex}
% ID: AIQ_L10C3_FULL_0500
% Mức: VD
Vật $A$ tác dụng lên vật $B$ lực $5\,\text{N}$. Theo định luật III Newton, độ lớn lực $B$ tác dụng lại $A$ bằng bao nhiêu?
\choiceTF
{\True Giá trị cần tìm là $5\,\text{N}$.}
{Giá trị cần tìm là $6\,\text{N}$.}
{\True Phải xác định đúng các lực tác dụng và chọn chiều dương trước khi lập phương trình.}
{Có thể bỏ qua điều kiện áp dụng và chiều của lực mà kết quả vẫn luôn đúng.}
\loigiai{
\begin{itemchoice}
\item (Đúng) Giá trị cần tìm là $5\,\text{N}$. (Vì): Hai lực tương tác có cùng độ lớn nên lực $B$ tác dụng lên $A$ bằng $5\,\text{N}$. b) Sai vì sai giá trị. c) Đúng. d) Sai vì phải kiểm tra điều kiện áp dụng và chiều lực
\item (Sai) Giá trị cần tìm là $6\,\text{N}$.
\item (Đúng) Phải xác định đúng các lực tác dụng và chọn chiều dương trước khi lập phương trình.
\item (Sai) Có thể bỏ qua điều kiện áp dụng và chiều của lực mà kết quả vẫn luôn đúng.
\end{itemchoice}
}
\end{ex}

% ===== Câu 15 =====
% ID: AIQ_L10C3_FULL_0501
% Mức: VD
\begin{ex}
% ID: AIQ_L10C3_FULL_0501
% Mức: VD
Vật $A$ tác dụng lên vật $B$ lực $7\,\text{N}$. Theo định luật III Newton, độ lớn lực $B$ tác dụng lại $A$ bằng bao nhiêu?
\choiceTF
{\True Giá trị cần tìm là $7\,\text{N}$.}
{Giá trị cần tìm là $8\,\text{N}$.}
{\True Phải xác định đúng các lực tác dụng và chọn chiều dương trước khi lập phương trình.}
{Có thể bỏ qua điều kiện áp dụng và chiều của lực mà kết quả vẫn luôn đúng.}
\loigiai{
\begin{itemchoice}
\item (Đúng) Giá trị cần tìm là $7\,\text{N}$. (Vì): Hai lực tương tác có cùng độ lớn nên lực $B$ tác dụng lên $A$ bằng $7\,\text{N}$. b) Sai vì sai giá trị. c) Đúng. d) Sai vì phải kiểm tra điều kiện áp dụng và chiều lực
\item (Sai) Giá trị cần tìm là $8\,\text{N}$.
\item (Đúng) Phải xác định đúng các lực tác dụng và chọn chiều dương trước khi lập phương trình.
\item (Sai) Có thể bỏ qua điều kiện áp dụng và chiều của lực mà kết quả vẫn luôn đúng.
\end{itemchoice}
}
\end{ex}

% ===== Câu 16 =====
% ID: AIQ_L10C3_FULL_0502
% Mức: TH
\begin{ex}
% ID: AIQ_L10C3_FULL_0502
% Mức: TH
Vật $A$ tác dụng lên vật $B$ lực $9\,\text{N}$. Theo định luật III Newton, độ lớn lực $B$ tác dụng lại $A$ bằng bao nhiêu?
\shortans{9}
\loigiai{
Hai lực tương tác có cùng độ lớn nên lực $B$ tác dụng lên $A$ bằng $9\,\text{N}$.
}
\end{ex}

% ===== Câu 17 =====
% ID: AIQ_L10C3_FULL_0503
% Mức: TH
\begin{ex}
% ID: AIQ_L10C3_FULL_0503
% Mức: TH
Vật $A$ tác dụng lên vật $B$ lực $11\,\text{N}$. Theo định luật III Newton, độ lớn lực $B$ tác dụng lại $A$ bằng bao nhiêu?
\shortans{11}
\loigiai{
Hai lực tương tác có cùng độ lớn nên lực $B$ tác dụng lên $A$ bằng $11\,\text{N}$.
}
\end{ex}

% ===== Câu 18 =====
% ID: AIQ_L10C3_FULL_0504
% Mức: VD
\begin{ex}
% ID: AIQ_L10C3_FULL_0504
% Mức: VD
Vật $A$ tác dụng lên vật $B$ lực $13\,\text{N}$. Theo định luật III Newton, độ lớn lực $B$ tác dụng lại $A$ bằng bao nhiêu?
\shortans{13}
\loigiai{
Hai lực tương tác có cùng độ lớn nên lực $B$ tác dụng lên $A$ bằng $13\,\text{N}$.
}
\end{ex}

% ===== Câu 19 =====
% ID: AIQ_L10C3_FULL_0505
% Mức: VD
\begin{ex}
% ID: AIQ_L10C3_FULL_0505
% Mức: VD
Vật $A$ tác dụng lên vật $B$ lực $15\,\text{N}$. Theo định luật III Newton, độ lớn lực $B$ tác dụng lại $A$ bằng bao nhiêu?
\shortans{15}
\loigiai{
Hai lực tương tác có cùng độ lớn nên lực $B$ tác dụng lên $A$ bằng $15\,\text{N}$.
}
\end{ex}

% ===== Câu 20 =====
% ID: AIQ_L10C3_FULL_0506
% Mức: VD
\begin{ex}
% ID: AIQ_L10C3_FULL_0506
% Mức: VD
Vật $A$ tác dụng lên vật $B$ lực $17\,\text{N}$. Theo định luật III Newton, độ lớn lực $B$ tác dụng lại $A$ bằng bao nhiêu?
\shortans{17}
\loigiai{
Hai lực tương tác có cùng độ lớn nên lực $B$ tác dụng lên $A$ bằng $17\,\text{N}$.
}
\end{ex}

% ===== Câu 21 =====
% ID: AIQ_L10C3_FULL_0507
% Mức: VDC
\begin{ex}
% ID: AIQ_L10C3_FULL_0507
% Mức: VDC
Vật $A$ tác dụng lên vật $B$ lực $19\,\text{N}$. Theo định luật III Newton, độ lớn lực $B$ tác dụng lại $A$ bằng bao nhiêu?
\shortans{19}
\loigiai{
Hai lực tương tác có cùng độ lớn nên lực $B$ tác dụng lên $A$ bằng $19\,\text{N}$.
}
\end{ex}

% ===== Câu 22 =====
% ID: AIQ_L10C3_FULL_0508
% Mức: VD
\begin{bt}
% ID: AIQ_L10C3_FULL_0508
% Mức: VD
Vật $A$ tác dụng lên vật $B$ lực $19\,\text{N}$. Theo định luật III Newton, độ lớn lực $B$ tác dụng lại $A$ bằng bao nhiêu?
\loigiai{
Phân tích lực, chọn hệ quy chiếu và áp dụng định luật hoặc quy tắc phù hợp. Hai lực tương tác có cùng độ lớn nên lực $B$ tác dụng lên $A$ bằng $19\,\text{N}$.
}
\end{bt}

% ===== Câu 23 =====
% ID: AIQ_L10C3_FULL_0509
% Mức: VD
\begin{bt}
% ID: AIQ_L10C3_FULL_0509
% Mức: VD
Vật $A$ tác dụng lên vật $B$ lực $5\,\text{N}$. Theo định luật III Newton, độ lớn lực $B$ tác dụng lại $A$ bằng bao nhiêu?
\loigiai{
Phân tích lực, chọn hệ quy chiếu và áp dụng định luật hoặc quy tắc phù hợp. Hai lực tương tác có cùng độ lớn nên lực $B$ tác dụng lên $A$ bằng $5\,\text{N}$.
}
\end{bt}

% ===== Câu 24 =====
% ID: AIQ_L10C3_FULL_0510
% Mức: VDC
\begin{bt}
% ID: AIQ_L10C3_FULL_0510
% Mức: VDC
Vật $A$ tác dụng lên vật $B$ lực $7\,\text{N}$. Theo định luật III Newton, độ lớn lực $B$ tác dụng lại $A$ bằng bao nhiêu?
\loigiai{
Phân tích lực, chọn hệ quy chiếu và áp dụng định luật hoặc quy tắc phù hợp. Hai lực tương tác có cùng độ lớn nên lực $B$ tác dụng lên $A$ bằng $7\,\text{N}$.
}
\end{bt}

% ===== Câu 25 =====
% ID: AIQ_L10C3_FULL_0511
% Mức: VDC
\begin{bt}
% ID: AIQ_L10C3_FULL_0511
% Mức: VDC
Vật $A$ tác dụng lên vật $B$ lực $9\,\text{N}$. Theo định luật III Newton, độ lớn lực $B$ tác dụng lại $A$ bằng bao nhiêu?
\loigiai{
Phân tích lực, chọn hệ quy chiếu và áp dụng định luật hoặc quy tắc phù hợp. Hai lực tương tác có cùng độ lớn nên lực $B$ tác dụng lên $A$ bằng $9\,\text{N}$.
}
\end{bt}

% ===== Câu 26 =====
% ID: AIQ_L10C3_FULL_0512
% Mức: VDC
\begin{bt}
% ID: AIQ_L10C3_FULL_0512
% Mức: VDC
Vật $A$ tác dụng lên vật $B$ lực $11\,\text{N}$. Theo định luật III Newton, độ lớn lực $B$ tác dụng lại $A$ bằng bao nhiêu?
\loigiai{
Phân tích lực, chọn hệ quy chiếu và áp dụng định luật hoặc quy tắc phù hợp. Hai lực tương tác có cùng độ lớn nên lực $B$ tác dụng lên $A$ bằng $11\,\text{N}$.
}
\end{bt}

% ===== Câu 27 =====
% ID: AIQ_L10C3_FULL_0513
% Mức: VDC
\begin{bt}
% ID: AIQ_L10C3_FULL_0513
% Mức: VDC
Vật $A$ tác dụng lên vật $B$ lực $13\,\text{N}$. Theo định luật III Newton, độ lớn lực $B$ tác dụng lại $A$ bằng bao nhiêu?
\loigiai{
Phân tích lực, chọn hệ quy chiếu và áp dụng định luật hoặc quy tắc phù hợp. Hai lực tương tác có cùng độ lớn nên lực $B$ tác dụng lên $A$ bằng $13\,\text{N}$.
}
\end{bt}

% ===== Câu 28 =====
% ID: AIQ_L10C3_FULL_0514
% Mức: NB
\dangbt{Phân biệt cặp lực trực đối với hai lực cân bằng}
\begin{ex}
% ID: AIQ_L10C3_FULL_0514
% Mức: NB
Hai vật $A$ và $B$ tương tác. Lực $A$ tác dụng $B$ có độ lớn $7\,\text{N}$. Xác định độ lớn của phản lực và nêu vật chịu từng lực.
\choice
{7 $N$}
{\True 7 $N$}
{6 $N$}
{14 $N$}
\loigiai{
Chọn B. Phản lực có độ lớn 7 $N$, đặt lên $A$; lực ban đầu đặt lên B. Vì đặt lên hai vật khác nhau nên chúng không phải hai lực cân bằng.
}
\end{ex}

% ===== Câu 29 =====
% ID: AIQ_L10C3_FULL_0515
% Mức: NB
\begin{ex}
% ID: AIQ_L10C3_FULL_0515
% Mức: NB
Hai vật $A$ và $B$ tương tác. Lực $A$ tác dụng $B$ có độ lớn $9\,\text{N}$. Xác định độ lớn của phản lực và nêu vật chịu từng lực.
\choice
{9 $N$}
{10 $N$}
{\True 9 $N$}
{18 $N$}
\loigiai{
Chọn C. Phản lực có độ lớn 9 $N$, đặt lên $A$; lực ban đầu đặt lên B. Vì đặt lên hai vật khác nhau nên chúng không phải hai lực cân bằng.
}
\end{ex}

% ===== Câu 30 =====
% ID: AIQ_L10C3_FULL_0516
% Mức: NB
\begin{ex}
% ID: AIQ_L10C3_FULL_0516
% Mức: NB
Hai vật $A$ và $B$ tương tác. Lực $A$ tác dụng $B$ có độ lớn $11\,\text{N}$. Xác định độ lớn của phản lực và nêu vật chịu từng lực.
\choice
{11 $N$}
{12 $N$}
{10 $N$}
{\True 11 $N$}
\loigiai{
Chọn D. Phản lực có độ lớn 11 $N$, đặt lên $A$; lực ban đầu đặt lên B. Vì đặt lên hai vật khác nhau nên chúng không phải hai lực cân bằng.
}
\end{ex}

% ===== Câu 31 =====
% ID: AIQ_L10C3_FULL_0517
% Mức: NB
\begin{ex}
% ID: AIQ_L10C3_FULL_0517
% Mức: NB
Hai vật $A$ và $B$ tương tác. Lực $A$ tác dụng $B$ có độ lớn $13\,\text{N}$. Xác định độ lớn của phản lực và nêu vật chịu từng lực.
\choice
{\True 13 $N$}
{14 $N$}
{12 $N$}
{26 $N$}
\loigiai{
Chọn A. Phản lực có độ lớn 13 $N$, đặt lên $A$; lực ban đầu đặt lên B. Vì đặt lên hai vật khác nhau nên chúng không phải hai lực cân bằng.
}
\end{ex}

% ===== Câu 32 =====
% ID: AIQ_L10C3_FULL_0518
% Mức: TH
\begin{ex}
% ID: AIQ_L10C3_FULL_0518
% Mức: TH
Hai vật $A$ và $B$ tương tác. Lực $A$ tác dụng $B$ có độ lớn $15\,\text{N}$. Xác định độ lớn của phản lực và nêu vật chịu từng lực.
\choice
{15 $N$}
{\True 15 $N$}
{14 $N$}
{30 $N$}
\loigiai{
Chọn B. Phản lực có độ lớn 15 $N$, đặt lên $A$; lực ban đầu đặt lên B. Vì đặt lên hai vật khác nhau nên chúng không phải hai lực cân bằng.
}
\end{ex}

% ===== Câu 33 =====
% ID: AIQ_L10C3_FULL_0519
% Mức: TH
\begin{ex}
% ID: AIQ_L10C3_FULL_0519
% Mức: TH
Hai vật $A$ và $B$ tương tác. Lực $A$ tác dụng $B$ có độ lớn $17\,\text{N}$. Xác định độ lớn của phản lực và nêu vật chịu từng lực.
\choice
{17 $N$}
{18 $N$}
{\True 17 $N$}
{34 $N$}
\loigiai{
Chọn C. Phản lực có độ lớn 17 $N$, đặt lên $A$; lực ban đầu đặt lên B. Vì đặt lên hai vật khác nhau nên chúng không phải hai lực cân bằng.
}
\end{ex}

% ===== Câu 34 =====
% ID: AIQ_L10C3_FULL_0520
% Mức: TH
\begin{ex}
% ID: AIQ_L10C3_FULL_0520
% Mức: TH
Hai vật $A$ và $B$ tương tác. Lực $A$ tác dụng $B$ có độ lớn $19\,\text{N}$. Xác định độ lớn của phản lực và nêu vật chịu từng lực.
\choice
{19 $N$}
{20 $N$}
{18 $N$}
{\True 19 $N$}
\loigiai{
Chọn D. Phản lực có độ lớn 19 $N$, đặt lên $A$; lực ban đầu đặt lên B. Vì đặt lên hai vật khác nhau nên chúng không phải hai lực cân bằng.
}
\end{ex}

% ===== Câu 35 =====
% ID: AIQ_L10C3_FULL_0521
% Mức: TH
\begin{ex}
% ID: AIQ_L10C3_FULL_0521
% Mức: TH
Hai vật $A$ và $B$ tương tác. Lực $A$ tác dụng $B$ có độ lớn $5\,\text{N}$. Xác định độ lớn của phản lực và nêu vật chịu từng lực.
\choice
{\True 5 $N$}
{6 $N$}
{4 $N$}
{10 $N$}
\loigiai{
Chọn A. Phản lực có độ lớn 5 $N$, đặt lên $A$; lực ban đầu đặt lên B. Vì đặt lên hai vật khác nhau nên chúng không phải hai lực cân bằng.
}
\end{ex}

% ===== Câu 36 =====
% ID: AIQ_L10C3_FULL_0522
% Mức: VD
\begin{ex}
% ID: AIQ_L10C3_FULL_0522
% Mức: VD
Hai vật $A$ và $B$ tương tác. Lực $A$ tác dụng $B$ có độ lớn $7\,\text{N}$. Xác định độ lớn của phản lực và nêu vật chịu từng lực.
\choice
{7 $N$}
{\True 7 $N$}
{6 $N$}
{14 $N$}
\loigiai{
Chọn B. Phản lực có độ lớn 7 $N$, đặt lên $A$; lực ban đầu đặt lên B. Vì đặt lên hai vật khác nhau nên chúng không phải hai lực cân bằng.
}
\end{ex}

% ===== Câu 37 =====
% ID: AIQ_L10C3_FULL_0523
% Mức: VD
\begin{ex}
% ID: AIQ_L10C3_FULL_0523
% Mức: VD
Hai vật $A$ và $B$ tương tác. Lực $A$ tác dụng $B$ có độ lớn $9\,\text{N}$. Xác định độ lớn của phản lực và nêu vật chịu từng lực.
\choice
{9 $N$}
{10 $N$}
{\True 9 $N$}
{18 $N$}
\loigiai{
Chọn C. Phản lực có độ lớn 9 $N$, đặt lên $A$; lực ban đầu đặt lên B. Vì đặt lên hai vật khác nhau nên chúng không phải hai lực cân bằng.
}
\end{ex}

% ===== Câu 38 =====
% ID: AIQ_L10C3_FULL_0524
% Mức: TH
\begin{ex}
% ID: AIQ_L10C3_FULL_0524
% Mức: TH
Hai vật $A$ và $B$ tương tác. Lực $A$ tác dụng $B$ có độ lớn $17\,\text{N}$. Xác định độ lớn của phản lực và nêu vật chịu từng lực.
\choiceTF
{\True Giá trị cần tìm là $17\,\text{N}$.}
{Giá trị cần tìm là $18\,\text{N}$.}
{\True Phải xác định đúng các lực tác dụng và chọn chiều dương trước khi lập phương trình.}
{Có thể bỏ qua điều kiện áp dụng và chiều của lực mà kết quả vẫn luôn đúng.}
\loigiai{
\begin{itemchoice}
\item (Đúng) Giá trị cần tìm là $17\,\text{N}$. (Vì): Phản lực có độ lớn 17 $N$, đặt lên $A$; lực ban đầu đặt lên B. Vì đặt lên hai vật khác nhau nên chúng không phải hai lực cân bằng. b) Sai vì sai giá trị. c) Đúng. d) Sai vì phải kiểm tra điều kiện áp dụng và chiều lực
\item (Sai) Giá trị cần tìm là $18\,\text{N}$.
\item (Đúng) Phải xác định đúng các lực tác dụng và chọn chiều dương trước khi lập phương trình.
\item (Sai) Có thể bỏ qua điều kiện áp dụng và chiều của lực mà kết quả vẫn luôn đúng.
\end{itemchoice}
}
\end{ex}

% ===== Câu 39 =====
% ID: AIQ_L10C3_FULL_0525
% Mức: TH
\begin{ex}
% ID: AIQ_L10C3_FULL_0525
% Mức: TH
Hai vật $A$ và $B$ tương tác. Lực $A$ tác dụng $B$ có độ lớn $19\,\text{N}$. Xác định độ lớn của phản lực và nêu vật chịu từng lực.
\choiceTF
{\True Giá trị cần tìm là $19\,\text{N}$.}
{Giá trị cần tìm là $20\,\text{N}$.}
{\True Phải xác định đúng các lực tác dụng và chọn chiều dương trước khi lập phương trình.}
{Có thể bỏ qua điều kiện áp dụng và chiều của lực mà kết quả vẫn luôn đúng.}
\loigiai{
\begin{itemchoice}
\item (Đúng) Giá trị cần tìm là $19\,\text{N}$. (Vì): Phản lực có độ lớn 19 $N$, đặt lên $A$; lực ban đầu đặt lên B. Vì đặt lên hai vật khác nhau nên chúng không phải hai lực cân bằng. b) Sai vì sai giá trị. c) Đúng. d) Sai vì phải kiểm tra điều kiện áp dụng và chiều lực
\item (Sai) Giá trị cần tìm là $20\,\text{N}$.
\item (Đúng) Phải xác định đúng các lực tác dụng và chọn chiều dương trước khi lập phương trình.
\item (Sai) Có thể bỏ qua điều kiện áp dụng và chiều của lực mà kết quả vẫn luôn đúng.
\end{itemchoice}
}
\end{ex}

% ===== Câu 40 =====
% ID: AIQ_L10C3_FULL_0526
% Mức: VD
\begin{ex}
% ID: AIQ_L10C3_FULL_0526
% Mức: VD
Hai vật $A$ và $B$ tương tác. Lực $A$ tác dụng $B$ có độ lớn $5\,\text{N}$. Xác định độ lớn của phản lực và nêu vật chịu từng lực.
\choiceTF
{\True Giá trị cần tìm là $5\,\text{N}$.}
{Giá trị cần tìm là $6\,\text{N}$.}
{\True Phải xác định đúng các lực tác dụng và chọn chiều dương trước khi lập phương trình.}
{Có thể bỏ qua điều kiện áp dụng và chiều của lực mà kết quả vẫn luôn đúng.}
\loigiai{
\begin{itemchoice}
\item (Đúng) Giá trị cần tìm là $5\,\text{N}$. (Vì): Phản lực có độ lớn 5 $N$, đặt lên $A$; lực ban đầu đặt lên B. Vì đặt lên hai vật khác nhau nên chúng không phải hai lực cân bằng. b) Sai vì sai giá trị. c) Đúng. d) Sai vì phải kiểm tra điều kiện áp dụng và chiều lực
\item (Sai) Giá trị cần tìm là $6\,\text{N}$.
\item (Đúng) Phải xác định đúng các lực tác dụng và chọn chiều dương trước khi lập phương trình.
\item (Sai) Có thể bỏ qua điều kiện áp dụng và chiều của lực mà kết quả vẫn luôn đúng.
\end{itemchoice}
}
\end{ex}

% ===== Câu 41 =====
% ID: AIQ_L10C3_FULL_0527
% Mức: VD
\begin{ex}
% ID: AIQ_L10C3_FULL_0527
% Mức: VD
Hai vật $A$ và $B$ tương tác. Lực $A$ tác dụng $B$ có độ lớn $7\,\text{N}$. Xác định độ lớn của phản lực và nêu vật chịu từng lực.
\choiceTF
{\True Giá trị cần tìm là $7\,\text{N}$.}
{Giá trị cần tìm là $8\,\text{N}$.}
{\True Phải xác định đúng các lực tác dụng và chọn chiều dương trước khi lập phương trình.}
{Có thể bỏ qua điều kiện áp dụng và chiều của lực mà kết quả vẫn luôn đúng.}
\loigiai{
\begin{itemchoice}
\item (Đúng) Giá trị cần tìm là $7\,\text{N}$. (Vì): Phản lực có độ lớn 7 $N$, đặt lên $A$; lực ban đầu đặt lên B. Vì đặt lên hai vật khác nhau nên chúng không phải hai lực cân bằng. b) Sai vì sai giá trị. c) Đúng. d) Sai vì phải kiểm tra điều kiện áp dụng và chiều lực
\item (Sai) Giá trị cần tìm là $8\,\text{N}$.
\item (Đúng) Phải xác định đúng các lực tác dụng và chọn chiều dương trước khi lập phương trình.
\item (Sai) Có thể bỏ qua điều kiện áp dụng và chiều của lực mà kết quả vẫn luôn đúng.
\end{itemchoice}
}
\end{ex}

% ===== Câu 42 =====
% ID: AIQ_L10C3_FULL_0528
% Mức: VD
\begin{ex}
% ID: AIQ_L10C3_FULL_0528
% Mức: VD
Hai vật $A$ và $B$ tương tác. Lực $A$ tác dụng $B$ có độ lớn $9\,\text{N}$. Xác định độ lớn của phản lực và nêu vật chịu từng lực.
\choiceTF
{\True Giá trị cần tìm là $9\,\text{N}$.}
{Giá trị cần tìm là $10\,\text{N}$.}
{\True Phải xác định đúng các lực tác dụng và chọn chiều dương trước khi lập phương trình.}
{Có thể bỏ qua điều kiện áp dụng và chiều của lực mà kết quả vẫn luôn đúng.}
\loigiai{
\begin{itemchoice}
\item (Đúng) Giá trị cần tìm là $9\,\text{N}$. (Vì): Phản lực có độ lớn 9 $N$, đặt lên $A$; lực ban đầu đặt lên B. Vì đặt lên hai vật khác nhau nên chúng không phải hai lực cân bằng. b) Sai vì sai giá trị. c) Đúng. d) Sai vì phải kiểm tra điều kiện áp dụng và chiều lực
\item (Sai) Giá trị cần tìm là $10\,\text{N}$.
\item (Đúng) Phải xác định đúng các lực tác dụng và chọn chiều dương trước khi lập phương trình.
\item (Sai) Có thể bỏ qua điều kiện áp dụng và chiều của lực mà kết quả vẫn luôn đúng.
\end{itemchoice}
}
\end{ex}

% ===== Câu 43 =====
% ID: AIQ_L10C3_FULL_0529
% Mức: TH
\begin{ex}
% ID: AIQ_L10C3_FULL_0529
% Mức: TH
Hai vật $A$ và $B$ tương tác. Lực $A$ tác dụng $B$ có độ lớn $11\,\text{N}$. Xác định độ lớn của phản lực và nêu vật chịu từng lực.
\shortans{11}
\loigiai{
Phản lực có độ lớn 11 $N$, đặt lên $A$; lực ban đầu đặt lên B. Vì đặt lên hai vật khác nhau nên chúng không phải hai lực cân bằng.
}
\end{ex}

% ===== Câu 44 =====
% ID: AIQ_L10C3_FULL_0530
% Mức: TH
\begin{ex}
% ID: AIQ_L10C3_FULL_0530
% Mức: TH
Hai vật $A$ và $B$ tương tác. Lực $A$ tác dụng $B$ có độ lớn $13\,\text{N}$. Xác định độ lớn của phản lực và nêu vật chịu từng lực.
\shortans{13}
\loigiai{
Phản lực có độ lớn 13 $N$, đặt lên $A$; lực ban đầu đặt lên B. Vì đặt lên hai vật khác nhau nên chúng không phải hai lực cân bằng.
}
\end{ex}

% ===== Câu 45 =====
% ID: AIQ_L10C3_FULL_0531
% Mức: VD
\begin{ex}
% ID: AIQ_L10C3_FULL_0531
% Mức: VD
Hai vật $A$ và $B$ tương tác. Lực $A$ tác dụng $B$ có độ lớn $15\,\text{N}$. Xác định độ lớn của phản lực và nêu vật chịu từng lực.
\shortans{15}
\loigiai{
Phản lực có độ lớn 15 $N$, đặt lên $A$; lực ban đầu đặt lên B. Vì đặt lên hai vật khác nhau nên chúng không phải hai lực cân bằng.
}
\end{ex}

% ===== Câu 46 =====
% ID: AIQ_L10C3_FULL_0532
% Mức: VD
\begin{ex}
% ID: AIQ_L10C3_FULL_0532
% Mức: VD
Hai vật $A$ và $B$ tương tác. Lực $A$ tác dụng $B$ có độ lớn $17\,\text{N}$. Xác định độ lớn của phản lực và nêu vật chịu từng lực.
\shortans{17}
\loigiai{
Phản lực có độ lớn 17 $N$, đặt lên $A$; lực ban đầu đặt lên B. Vì đặt lên hai vật khác nhau nên chúng không phải hai lực cân bằng.
}
\end{ex}

% ===== Câu 47 =====
% ID: AIQ_L10C3_FULL_0533
% Mức: VD
\begin{ex}
% ID: AIQ_L10C3_FULL_0533
% Mức: VD
Hai vật $A$ và $B$ tương tác. Lực $A$ tác dụng $B$ có độ lớn $19\,\text{N}$. Xác định độ lớn của phản lực và nêu vật chịu từng lực.
\shortans{19}
\loigiai{
Phản lực có độ lớn 19 $N$, đặt lên $A$; lực ban đầu đặt lên B. Vì đặt lên hai vật khác nhau nên chúng không phải hai lực cân bằng.
}
\end{ex}

% ===== Câu 48 =====
% ID: AIQ_L10C3_FULL_0534
% Mức: VDC
\begin{ex}
% ID: AIQ_L10C3_FULL_0534
% Mức: VDC
Hai vật $A$ và $B$ tương tác. Lực $A$ tác dụng $B$ có độ lớn $5\,\text{N}$. Xác định độ lớn của phản lực và nêu vật chịu từng lực.
\shortans{5}
\loigiai{
Phản lực có độ lớn 5 $N$, đặt lên $A$; lực ban đầu đặt lên B. Vì đặt lên hai vật khác nhau nên chúng không phải hai lực cân bằng.
}
\end{ex}

% ===== Câu 49 =====
% ID: AIQ_L10C3_FULL_0535
% Mức: VD
\begin{bt}
% ID: AIQ_L10C3_FULL_0535
% Mức: VD
Hai vật $A$ và $B$ tương tác. Lực $A$ tác dụng $B$ có độ lớn $5\,\text{N}$. Xác định độ lớn của phản lực và nêu vật chịu từng lực.
\loigiai{
Phân tích lực, chọn hệ quy chiếu và áp dụng định luật hoặc quy tắc phù hợp. Phản lực có độ lớn 5 $N$, đặt lên $A$; lực ban đầu đặt lên B. Vì đặt lên hai vật khác nhau nên chúng không phải hai lực cân bằng.
}
\end{bt}

% ===== Câu 50 =====
% ID: AIQ_L10C3_FULL_0536
% Mức: VD
\begin{bt}
% ID: AIQ_L10C3_FULL_0536
% Mức: VD
Hai vật $A$ và $B$ tương tác. Lực $A$ tác dụng $B$ có độ lớn $7\,\text{N}$. Xác định độ lớn của phản lực và nêu vật chịu từng lực.
\loigiai{
Phân tích lực, chọn hệ quy chiếu và áp dụng định luật hoặc quy tắc phù hợp. Phản lực có độ lớn 7 $N$, đặt lên $A$; lực ban đầu đặt lên B. Vì đặt lên hai vật khác nhau nên chúng không phải hai lực cân bằng.
}
\end{bt}

% ===== Câu 51 =====
% ID: AIQ_L10C3_FULL_0537
% Mức: VDC
\begin{bt}
% ID: AIQ_L10C3_FULL_0537
% Mức: VDC
Hai vật $A$ và $B$ tương tác. Lực $A$ tác dụng $B$ có độ lớn $9\,\text{N}$. Xác định độ lớn của phản lực và nêu vật chịu từng lực.
\loigiai{
Phân tích lực, chọn hệ quy chiếu và áp dụng định luật hoặc quy tắc phù hợp. Phản lực có độ lớn 9 $N$, đặt lên $A$; lực ban đầu đặt lên B. Vì đặt lên hai vật khác nhau nên chúng không phải hai lực cân bằng.
}
\end{bt}

% ===== Câu 52 =====
% ID: AIQ_L10C3_FULL_0538
% Mức: VDC
\begin{bt}
% ID: AIQ_L10C3_FULL_0538
% Mức: VDC
Hai vật $A$ và $B$ tương tác. Lực $A$ tác dụng $B$ có độ lớn $11\,\text{N}$. Xác định độ lớn của phản lực và nêu vật chịu từng lực.
\loigiai{
Phân tích lực, chọn hệ quy chiếu và áp dụng định luật hoặc quy tắc phù hợp. Phản lực có độ lớn 11 $N$, đặt lên $A$; lực ban đầu đặt lên B. Vì đặt lên hai vật khác nhau nên chúng không phải hai lực cân bằng.
}
\end{bt}

% ===== Câu 53 =====
% ID: AIQ_L10C3_FULL_0539
% Mức: VDC
\begin{bt}
% ID: AIQ_L10C3_FULL_0539
% Mức: VDC
Hai vật $A$ và $B$ tương tác. Lực $A$ tác dụng $B$ có độ lớn $13\,\text{N}$. Xác định độ lớn của phản lực và nêu vật chịu từng lực.
\loigiai{
Phân tích lực, chọn hệ quy chiếu và áp dụng định luật hoặc quy tắc phù hợp. Phản lực có độ lớn 13 $N$, đặt lên $A$; lực ban đầu đặt lên B. Vì đặt lên hai vật khác nhau nên chúng không phải hai lực cân bằng.
}
\end{bt}

% ===== Câu 54 =====
% ID: AIQ_L10C3_FULL_0540
% Mức: VDC
\begin{bt}
% ID: AIQ_L10C3_FULL_0540
% Mức: VDC
Hai vật $A$ và $B$ tương tác. Lực $A$ tác dụng $B$ có độ lớn $15\,\text{N}$. Xác định độ lớn của phản lực và nêu vật chịu từng lực.
\loigiai{
Phân tích lực, chọn hệ quy chiếu và áp dụng định luật hoặc quy tắc phù hợp. Phản lực có độ lớn 15 $N$, đặt lên $A$; lực ban đầu đặt lên B. Vì đặt lên hai vật khác nhau nên chúng không phải hai lực cân bằng.
}
\end{bt}

% ===== Câu 55 =====
% ID: AIQ_L10C3_FULL_0541
% Mức: NB
\dangbt{Xác định phương, chiều và độ lớn của lực tương tác}
\begin{ex}
% ID: AIQ_L10C3_FULL_0541
% Mức: NB
Vật $A$ đẩy vật $B$ theo phương ngang sang phải với lực $9\,\text{N}$. Chọn chiều dương sang phải. Giá trị đại số của lực $B$ tác dụng lên $A$ là bao nhiêu?
\choice
{-9 $N$}
{-8 $N$}
{\True -9 $N$}
{-18 $N$}
\loigiai{
Chọn C. Lực phản lực ngược chiều lực $A$ tác dụng $B$ nên có giá trị đại số − $9\,\text{N}$.
}
\end{ex}

% ===== Câu 56 =====
% ID: AIQ_L10C3_FULL_0542
% Mức: NB
\begin{ex}
% ID: AIQ_L10C3_FULL_0542
% Mức: NB
Vật $A$ đẩy vật $B$ theo phương ngang sang phải với lực $11\,\text{N}$. Chọn chiều dương sang phải. Giá trị đại số của lực $B$ tác dụng lên $A$ là bao nhiêu?
\choice
{-11 $N$}
{-10 $N$}
{0 $N$}
{\True -11 $N$}
\loigiai{
Chọn D. Lực phản lực ngược chiều lực $A$ tác dụng $B$ nên có giá trị đại số − $11\,\text{N}$.
}
\end{ex}

% ===== Câu 57 =====
% ID: AIQ_L10C3_FULL_0543
% Mức: NB
\begin{ex}
% ID: AIQ_L10C3_FULL_0543
% Mức: NB
Vật $A$ đẩy vật $B$ theo phương ngang sang phải với lực $13\,\text{N}$. Chọn chiều dương sang phải. Giá trị đại số của lực $B$ tác dụng lên $A$ là bao nhiêu?
\choice
{\True -13 $N$}
{-12 $N$}
{0 $N$}
{-26 $N$}
\loigiai{
Chọn A. Lực phản lực ngược chiều lực $A$ tác dụng $B$ nên có giá trị đại số − $13\,\text{N}$.
}
\end{ex}

% ===== Câu 58 =====
% ID: AIQ_L10C3_FULL_0544
% Mức: NB
\begin{ex}
% ID: AIQ_L10C3_FULL_0544
% Mức: NB
Vật $A$ đẩy vật $B$ theo phương ngang sang phải với lực $15\,\text{N}$. Chọn chiều dương sang phải. Giá trị đại số của lực $B$ tác dụng lên $A$ là bao nhiêu?
\choice
{-15 $N$}
{\True -15 $N$}
{0 $N$}
{-30 $N$}
\loigiai{
Chọn B. Lực phản lực ngược chiều lực $A$ tác dụng $B$ nên có giá trị đại số − $15\,\text{N}$.
}
\end{ex}

% ===== Câu 59 =====
% ID: AIQ_L10C3_FULL_0545
% Mức: TH
\begin{ex}
% ID: AIQ_L10C3_FULL_0545
% Mức: TH
Vật $A$ đẩy vật $B$ theo phương ngang sang phải với lực $17\,\text{N}$. Chọn chiều dương sang phải. Giá trị đại số của lực $B$ tác dụng lên $A$ là bao nhiêu?
\choice
{-17 $N$}
{-16 $N$}
{\True -17 $N$}
{-34 $N$}
\loigiai{
Chọn C. Lực phản lực ngược chiều lực $A$ tác dụng $B$ nên có giá trị đại số − $17\,\text{N}$.
}
\end{ex}

% ===== Câu 60 =====
% ID: AIQ_L10C3_FULL_0546
% Mức: TH
\begin{ex}
% ID: AIQ_L10C3_FULL_0546
% Mức: TH
Vật $A$ đẩy vật $B$ theo phương ngang sang phải với lực $19\,\text{N}$. Chọn chiều dương sang phải. Giá trị đại số của lực $B$ tác dụng lên $A$ là bao nhiêu?
\choice
{-19 $N$}
{-18 $N$}
{0 $N$}
{\True -19 $N$}
\loigiai{
Chọn D. Lực phản lực ngược chiều lực $A$ tác dụng $B$ nên có giá trị đại số − $19\,\text{N}$.
}
\end{ex}

% ===== Câu 61 =====
% ID: AIQ_L10C3_FULL_0547
% Mức: TH
\begin{ex}
% ID: AIQ_L10C3_FULL_0547
% Mức: TH
Vật $A$ đẩy vật $B$ theo phương ngang sang phải với lực $5\,\text{N}$. Chọn chiều dương sang phải. Giá trị đại số của lực $B$ tác dụng lên $A$ là bao nhiêu?
\choice
{\True -5 $N$}
{-4 $N$}
{0 $N$}
{-10 $N$}
\loigiai{
Chọn A. Lực phản lực ngược chiều lực $A$ tác dụng $B$ nên có giá trị đại số − $5\,\text{N}$.
}
\end{ex}

% ===== Câu 62 =====
% ID: AIQ_L10C3_FULL_0548
% Mức: TH
\begin{ex}
% ID: AIQ_L10C3_FULL_0548
% Mức: TH
Vật $A$ đẩy vật $B$ theo phương ngang sang phải với lực $7\,\text{N}$. Chọn chiều dương sang phải. Giá trị đại số của lực $B$ tác dụng lên $A$ là bao nhiêu?
\choice
{-7 $N$}
{\True -7 $N$}
{0 $N$}
{-14 $N$}
\loigiai{
Chọn B. Lực phản lực ngược chiều lực $A$ tác dụng $B$ nên có giá trị đại số − $7\,\text{N}$.
}
\end{ex}

% ===== Câu 63 =====
% ID: AIQ_L10C3_FULL_0549
% Mức: VD
\begin{ex}
% ID: AIQ_L10C3_FULL_0549
% Mức: VD
Vật $A$ đẩy vật $B$ theo phương ngang sang phải với lực $9\,\text{N}$. Chọn chiều dương sang phải. Giá trị đại số của lực $B$ tác dụng lên $A$ là bao nhiêu?
\choice
{-9 $N$}
{-8 $N$}
{\True -9 $N$}
{-18 $N$}
\loigiai{
Chọn C. Lực phản lực ngược chiều lực $A$ tác dụng $B$ nên có giá trị đại số − $9\,\text{N}$.
}
\end{ex}

% ===== Câu 64 =====
% ID: AIQ_L10C3_FULL_0550
% Mức: VD
\begin{ex}
% ID: AIQ_L10C3_FULL_0550
% Mức: VD
Vật $A$ đẩy vật $B$ theo phương ngang sang phải với lực $11\,\text{N}$. Chọn chiều dương sang phải. Giá trị đại số của lực $B$ tác dụng lên $A$ là bao nhiêu?
\choice
{-11 $N$}
{-10 $N$}
{0 $N$}
{\True -11 $N$}
\loigiai{
Chọn D. Lực phản lực ngược chiều lực $A$ tác dụng $B$ nên có giá trị đại số − $11\,\text{N}$.
}
\end{ex}

% ===== Câu 65 =====
% ID: AIQ_L10C3_FULL_0551
% Mức: TH
\begin{ex}
% ID: AIQ_L10C3_FULL_0551
% Mức: TH
Vật $A$ đẩy vật $B$ theo phương ngang sang phải với lực $19\,\text{N}$. Chọn chiều dương sang phải. Giá trị đại số của lực $B$ tác dụng lên $A$ là bao nhiêu?
\choiceTF
{\True Giá trị cần tìm là - $19\,\text{N}$.}
{Giá trị cần tìm là - $18\,\text{N}$.}
{\True Phải xác định đúng các lực tác dụng và chọn chiều dương trước khi lập phương trình.}
{Có thể bỏ qua điều kiện áp dụng và chiều của lực mà kết quả vẫn luôn đúng.}
\loigiai{
\begin{itemchoice}
\item (Đúng) Giá trị cần tìm là - $19\,\text{N}$. (Vì): Lực phản lực ngược chiều lực $A$ tác dụng $B$ nên có giá trị đại số − $19\,\text{N}$. b) Sai vì sai giá trị. c) Đúng. d) Sai vì phải kiểm tra điều kiện áp dụng và chiều lực
\item (Sai) Giá trị cần tìm là - $18\,\text{N}$.
\item (Đúng) Phải xác định đúng các lực tác dụng và chọn chiều dương trước khi lập phương trình.
\item (Sai) Có thể bỏ qua điều kiện áp dụng và chiều của lực mà kết quả vẫn luôn đúng.
\end{itemchoice}
}
\end{ex}

% ===== Câu 66 =====
% ID: AIQ_L10C3_FULL_0552
% Mức: TH
\begin{ex}
% ID: AIQ_L10C3_FULL_0552
% Mức: TH
Vật $A$ đẩy vật $B$ theo phương ngang sang phải với lực $5\,\text{N}$. Chọn chiều dương sang phải. Giá trị đại số của lực $B$ tác dụng lên $A$ là bao nhiêu?
\choiceTF
{\True Giá trị cần tìm là - $5\,\text{N}$.}
{Giá trị cần tìm là - $4\,\text{N}$.}
{\True Phải xác định đúng các lực tác dụng và chọn chiều dương trước khi lập phương trình.}
{Có thể bỏ qua điều kiện áp dụng và chiều của lực mà kết quả vẫn luôn đúng.}
\loigiai{
\begin{itemchoice}
\item (Đúng) Giá trị cần tìm là - $5\,\text{N}$. (Vì): Lực phản lực ngược chiều lực $A$ tác dụng $B$ nên có giá trị đại số − $5\,\text{N}$. b) Sai vì sai giá trị. c) Đúng. d) Sai vì phải kiểm tra điều kiện áp dụng và chiều lực
\item (Sai) Giá trị cần tìm là - $4\,\text{N}$.
\item (Đúng) Phải xác định đúng các lực tác dụng và chọn chiều dương trước khi lập phương trình.
\item (Sai) Có thể bỏ qua điều kiện áp dụng và chiều của lực mà kết quả vẫn luôn đúng.
\end{itemchoice}
}
\end{ex}

% ===== Câu 67 =====
% ID: AIQ_L10C3_FULL_0553
% Mức: VD
\begin{ex}
% ID: AIQ_L10C3_FULL_0553
% Mức: VD
Vật $A$ đẩy vật $B$ theo phương ngang sang phải với lực $7\,\text{N}$. Chọn chiều dương sang phải. Giá trị đại số của lực $B$ tác dụng lên $A$ là bao nhiêu?
\choiceTF
{\True Giá trị cần tìm là - $7\,\text{N}$.}
{Giá trị cần tìm là - $6\,\text{N}$.}
{\True Phải xác định đúng các lực tác dụng và chọn chiều dương trước khi lập phương trình.}
{Có thể bỏ qua điều kiện áp dụng và chiều của lực mà kết quả vẫn luôn đúng.}
\loigiai{
\begin{itemchoice}
\item (Đúng) Giá trị cần tìm là - $7\,\text{N}$. (Vì): Lực phản lực ngược chiều lực $A$ tác dụng $B$ nên có giá trị đại số − $7\,\text{N}$. b) Sai vì sai giá trị. c) Đúng. d) Sai vì phải kiểm tra điều kiện áp dụng và chiều lực
\item (Sai) Giá trị cần tìm là - $6\,\text{N}$.
\item (Đúng) Phải xác định đúng các lực tác dụng và chọn chiều dương trước khi lập phương trình.
\item (Sai) Có thể bỏ qua điều kiện áp dụng và chiều của lực mà kết quả vẫn luôn đúng.
\end{itemchoice}
}
\end{ex}

% ===== Câu 68 =====
% ID: AIQ_L10C3_FULL_0554
% Mức: VD
\begin{ex}
% ID: AIQ_L10C3_FULL_0554
% Mức: VD
Vật $A$ đẩy vật $B$ theo phương ngang sang phải với lực $9\,\text{N}$. Chọn chiều dương sang phải. Giá trị đại số của lực $B$ tác dụng lên $A$ là bao nhiêu?
\choiceTF
{\True Giá trị cần tìm là - $9\,\text{N}$.}
{Giá trị cần tìm là - $8\,\text{N}$.}
{\True Phải xác định đúng các lực tác dụng và chọn chiều dương trước khi lập phương trình.}
{Có thể bỏ qua điều kiện áp dụng và chiều của lực mà kết quả vẫn luôn đúng.}
\loigiai{
\begin{itemchoice}
\item (Đúng) Giá trị cần tìm là - $9\,\text{N}$. (Vì): Lực phản lực ngược chiều lực $A$ tác dụng $B$ nên có giá trị đại số − $9\,\text{N}$. b) Sai vì sai giá trị. c) Đúng. d) Sai vì phải kiểm tra điều kiện áp dụng và chiều lực
\item (Sai) Giá trị cần tìm là - $8\,\text{N}$.
\item (Đúng) Phải xác định đúng các lực tác dụng và chọn chiều dương trước khi lập phương trình.
\item (Sai) Có thể bỏ qua điều kiện áp dụng và chiều của lực mà kết quả vẫn luôn đúng.
\end{itemchoice}
}
\end{ex}

% ===== Câu 69 =====
% ID: AIQ_L10C3_FULL_0555
% Mức: VD
\begin{ex}
% ID: AIQ_L10C3_FULL_0555
% Mức: VD
Vật $A$ đẩy vật $B$ theo phương ngang sang phải với lực $11\,\text{N}$. Chọn chiều dương sang phải. Giá trị đại số của lực $B$ tác dụng lên $A$ là bao nhiêu?
\choiceTF
{\True Giá trị cần tìm là - $11\,\text{N}$.}
{Giá trị cần tìm là - $10\,\text{N}$.}
{\True Phải xác định đúng các lực tác dụng và chọn chiều dương trước khi lập phương trình.}
{Có thể bỏ qua điều kiện áp dụng và chiều của lực mà kết quả vẫn luôn đúng.}
\loigiai{
\begin{itemchoice}
\item (Đúng) Giá trị cần tìm là - $11\,\text{N}$. (Vì): Lực phản lực ngược chiều lực $A$ tác dụng $B$ nên có giá trị đại số − $11\,\text{N}$. b) Sai vì sai giá trị. c) Đúng. d) Sai vì phải kiểm tra điều kiện áp dụng và chiều lực
\item (Sai) Giá trị cần tìm là - $10\,\text{N}$.
\item (Đúng) Phải xác định đúng các lực tác dụng và chọn chiều dương trước khi lập phương trình.
\item (Sai) Có thể bỏ qua điều kiện áp dụng và chiều của lực mà kết quả vẫn luôn đúng.
\end{itemchoice}
}
\end{ex}

% ===== Câu 70 =====
% ID: AIQ_L10C3_FULL_0556
% Mức: TH
\begin{ex}
% ID: AIQ_L10C3_FULL_0556
% Mức: TH
Vật $A$ đẩy vật $B$ theo phương ngang sang phải với lực $13\,\text{N}$. Chọn chiều dương sang phải. Giá trị đại số của lực $B$ tác dụng lên $A$ là bao nhiêu?
\shortans{-13}
\loigiai{
Lực phản lực ngược chiều lực $A$ tác dụng $B$ nên có giá trị đại số − $13\,\text{N}$.
}
\end{ex}

% ===== Câu 71 =====
% ID: AIQ_L10C3_FULL_0557
% Mức: TH
\begin{ex}
% ID: AIQ_L10C3_FULL_0557
% Mức: TH
Vật $A$ đẩy vật $B$ theo phương ngang sang phải với lực $15\,\text{N}$. Chọn chiều dương sang phải. Giá trị đại số của lực $B$ tác dụng lên $A$ là bao nhiêu?
\shortans{-15}
\loigiai{
Lực phản lực ngược chiều lực $A$ tác dụng $B$ nên có giá trị đại số − $15\,\text{N}$.
}
\end{ex}

% ===== Câu 72 =====
% ID: AIQ_L10C3_FULL_0558
% Mức: VD
\begin{ex}
% ID: AIQ_L10C3_FULL_0558
% Mức: VD
Vật $A$ đẩy vật $B$ theo phương ngang sang phải với lực $17\,\text{N}$. Chọn chiều dương sang phải. Giá trị đại số của lực $B$ tác dụng lên $A$ là bao nhiêu?
\shortans{-17}
\loigiai{
Lực phản lực ngược chiều lực $A$ tác dụng $B$ nên có giá trị đại số − $17\,\text{N}$.
}
\end{ex}

% ===== Câu 73 =====
% ID: AIQ_L10C3_FULL_0559
% Mức: VD
\begin{ex}
% ID: AIQ_L10C3_FULL_0559
% Mức: VD
Vật $A$ đẩy vật $B$ theo phương ngang sang phải với lực $19\,\text{N}$. Chọn chiều dương sang phải. Giá trị đại số của lực $B$ tác dụng lên $A$ là bao nhiêu?
\shortans{-19}
\loigiai{
Lực phản lực ngược chiều lực $A$ tác dụng $B$ nên có giá trị đại số − $19\,\text{N}$.
}
\end{ex}

% ===== Câu 74 =====
% ID: AIQ_L10C3_FULL_0560
% Mức: VD
\begin{ex}
% ID: AIQ_L10C3_FULL_0560
% Mức: VD
Vật $A$ đẩy vật $B$ theo phương ngang sang phải với lực $5\,\text{N}$. Chọn chiều dương sang phải. Giá trị đại số của lực $B$ tác dụng lên $A$ là bao nhiêu?
\shortans{-5}
\loigiai{
Lực phản lực ngược chiều lực $A$ tác dụng $B$ nên có giá trị đại số − $5\,\text{N}$.
}
\end{ex}

% ===== Câu 75 =====
% ID: AIQ_L10C3_FULL_0561
% Mức: VDC
\begin{ex}
% ID: AIQ_L10C3_FULL_0561
% Mức: VDC
Vật $A$ đẩy vật $B$ theo phương ngang sang phải với lực $7\,\text{N}$. Chọn chiều dương sang phải. Giá trị đại số của lực $B$ tác dụng lên $A$ là bao nhiêu?
\shortans{-7}
\loigiai{
Lực phản lực ngược chiều lực $A$ tác dụng $B$ nên có giá trị đại số − $7\,\text{N}$.
}
\end{ex}

% ===== Câu 76 =====
% ID: AIQ_L10C3_FULL_0562
% Mức: VD
\begin{bt}
% ID: AIQ_L10C3_FULL_0562
% Mức: VD
Vật $A$ đẩy vật $B$ theo phương ngang sang phải với lực $7\,\text{N}$. Chọn chiều dương sang phải. Giá trị đại số của lực $B$ tác dụng lên $A$ là bao nhiêu?
\loigiai{
Phân tích lực, chọn hệ quy chiếu và áp dụng định luật hoặc quy tắc phù hợp. Lực phản lực ngược chiều lực $A$ tác dụng $B$ nên có giá trị đại số − $7\,\text{N}$.
}
\end{bt}

% ===== Câu 77 =====
% ID: AIQ_L10C3_FULL_0563
% Mức: VD
\begin{bt}
% ID: AIQ_L10C3_FULL_0563
% Mức: VD
Vật $A$ đẩy vật $B$ theo phương ngang sang phải với lực $9\,\text{N}$. Chọn chiều dương sang phải. Giá trị đại số của lực $B$ tác dụng lên $A$ là bao nhiêu?
\loigiai{
Phân tích lực, chọn hệ quy chiếu và áp dụng định luật hoặc quy tắc phù hợp. Lực phản lực ngược chiều lực $A$ tác dụng $B$ nên có giá trị đại số − $9\,\text{N}$.
}
\end{bt}

% ===== Câu 78 =====
% ID: AIQ_L10C3_FULL_0564
% Mức: VDC
\begin{bt}
% ID: AIQ_L10C3_FULL_0564
% Mức: VDC
Vật $A$ đẩy vật $B$ theo phương ngang sang phải với lực $11\,\text{N}$. Chọn chiều dương sang phải. Giá trị đại số của lực $B$ tác dụng lên $A$ là bao nhiêu?
\loigiai{
Phân tích lực, chọn hệ quy chiếu và áp dụng định luật hoặc quy tắc phù hợp. Lực phản lực ngược chiều lực $A$ tác dụng $B$ nên có giá trị đại số − $11\,\text{N}$.
}
\end{bt}

% ===== Câu 79 =====
% ID: AIQ_L10C3_FULL_0565
% Mức: VDC
\begin{bt}
% ID: AIQ_L10C3_FULL_0565
% Mức: VDC
Vật $A$ đẩy vật $B$ theo phương ngang sang phải với lực $13\,\text{N}$. Chọn chiều dương sang phải. Giá trị đại số của lực $B$ tác dụng lên $A$ là bao nhiêu?
\loigiai{
Phân tích lực, chọn hệ quy chiếu và áp dụng định luật hoặc quy tắc phù hợp. Lực phản lực ngược chiều lực $A$ tác dụng $B$ nên có giá trị đại số − $13\,\text{N}$.
}
\end{bt}

% ===== Câu 80 =====
% ID: AIQ_L10C3_FULL_0566
% Mức: VDC
\begin{bt}
% ID: AIQ_L10C3_FULL_0566
% Mức: VDC
Vật $A$ đẩy vật $B$ theo phương ngang sang phải với lực $15\,\text{N}$. Chọn chiều dương sang phải. Giá trị đại số của lực $B$ tác dụng lên $A$ là bao nhiêu?
\loigiai{
Phân tích lực, chọn hệ quy chiếu và áp dụng định luật hoặc quy tắc phù hợp. Lực phản lực ngược chiều lực $A$ tác dụng $B$ nên có giá trị đại số − $15\,\text{N}$.
}
\end{bt}

% ===== Câu 81 =====
% ID: AIQ_L10C3_FULL_0567
% Mức: VDC
\begin{bt}
% ID: AIQ_L10C3_FULL_0567
% Mức: VDC
Vật $A$ đẩy vật $B$ theo phương ngang sang phải với lực $17\,\text{N}$. Chọn chiều dương sang phải. Giá trị đại số của lực $B$ tác dụng lên $A$ là bao nhiêu?
\loigiai{
Phân tích lực, chọn hệ quy chiếu và áp dụng định luật hoặc quy tắc phù hợp. Lực phản lực ngược chiều lực $A$ tác dụng $B$ nên có giá trị đại số − $17\,\text{N}$.
}
\end{bt}

% ===== Câu 82 =====
% ID: AIQ_L10C3_FULL_0568
% Mức: NB
\dangbt{Vận dụng định luật III Newton trong va chạm và tương tác}
\begin{ex}
% ID: AIQ_L10C3_FULL_0568
% Mức: NB
Hai xe va chạm. Xe 1 chịu lực trung bình 11 $N$ từ xe 2. Khối lượng xe 1 là $2\,\text{kg}$. Tính độ lớn gia tốc trung bình của xe 1.
\choice
{$5,5 m/s^2$}
{$6,5 m/s^2$}
{$4,5 m/s^2$}
{\True $5,5 m/s^2$}
\loigiai{
Chọn D. Theo định luật II Newton, a=F/m=11/2= $5,5\,\text{m}$ /s^2; lực tương tác trên xe 2 có cùng độ lớn.
}
\end{ex}

% ===== Câu 83 =====
% ID: AIQ_L10C3_FULL_0569
% Mức: NB
\begin{ex}
% ID: AIQ_L10C3_FULL_0569
% Mức: NB
Hai xe va chạm. Xe 1 chịu lực trung bình 13 $N$ từ xe 2. Khối lượng xe 1 là $3\,\text{kg}$. Tính độ lớn gia tốc trung bình của xe 1.
\choice
{\True $4,33 m/s^2$}
{$5,33 m/s^2$}
{$3,33 m/s^2$}
{$8,66 m/s^2$}
\loigiai{
Chọn A. Theo định luật II Newton, a=F/m=13/3= $4,33\,\text{m}$ /s^2; lực tương tác trên xe 2 có cùng độ lớn.
}
\end{ex}

% ===== Câu 84 =====
% ID: AIQ_L10C3_FULL_0570
% Mức: NB
\begin{ex}
% ID: AIQ_L10C3_FULL_0570
% Mức: NB
Hai xe va chạm. Xe 1 chịu lực trung bình 15 $N$ từ xe 2. Khối lượng xe 1 là $4\,\text{kg}$. Tính độ lớn gia tốc trung bình của xe 1.
\choice
{$3,75 m/s^2$}
{\True $3,75 m/s^2$}
{$2,75 m/s^2$}
{$7,5 m/s^2$}
\loigiai{
Chọn B. Theo định luật II Newton, a=F/m=15/4= $3,75\,\text{m}$ /s^2; lực tương tác trên xe 2 có cùng độ lớn.
}
\end{ex}

% ===== Câu 85 =====
% ID: AIQ_L10C3_FULL_0571
% Mức: NB
\begin{ex}
% ID: AIQ_L10C3_FULL_0571
% Mức: NB
Hai xe va chạm. Xe 1 chịu lực trung bình 17 $N$ từ xe 2. Khối lượng xe 1 là $5\,\text{kg}$. Tính độ lớn gia tốc trung bình của xe 1.
\choice
{$3,4 m/s^2$}
{$4,4 m/s^2$}
{\True $3,4 m/s^2$}
{$6,8 m/s^2$}
\loigiai{
Chọn C. Theo định luật II Newton, a=F/m=17/5= $3,4\,\text{m}$ /s^2; lực tương tác trên xe 2 có cùng độ lớn.
}
\end{ex}

% ===== Câu 86 =====
% ID: AIQ_L10C3_FULL_0572
% Mức: TH
\begin{ex}
% ID: AIQ_L10C3_FULL_0572
% Mức: TH
Hai xe va chạm. Xe 1 chịu lực trung bình 19 $N$ từ xe 2. Khối lượng xe 1 là $6\,\text{kg}$. Tính độ lớn gia tốc trung bình của xe 1.
\choice
{$3,17 m/s^2$}
{$4,17 m/s^2$}
{$2,17 m/s^2$}
{\True $3,17 m/s^2$}
\loigiai{
Chọn D. Theo định luật II Newton, a=F/m=19/6= $3,17\,\text{m}$ /s^2; lực tương tác trên xe 2 có cùng độ lớn.
}
\end{ex}

% ===== Câu 87 =====
% ID: AIQ_L10C3_FULL_0573
% Mức: TH
\begin{ex}
% ID: AIQ_L10C3_FULL_0573
% Mức: TH
Hai xe va chạm. Xe 1 chịu lực trung bình 5 $N$ từ xe 2. Khối lượng xe 1 là $2\,\text{kg}$. Tính độ lớn gia tốc trung bình của xe 1.
\choice
{\True $2,5 m/s^2$}
{$3,5 m/s^2$}
{$1,5 m/s^2$}
{$5 m/s^2$}
\loigiai{
Chọn A. Theo định luật II Newton, a=F/m=5/2= $2,5\,\text{m}$ /s^2; lực tương tác trên xe 2 có cùng độ lớn.
}
\end{ex}

% ===== Câu 88 =====
% ID: AIQ_L10C3_FULL_0574
% Mức: TH
\begin{ex}
% ID: AIQ_L10C3_FULL_0574
% Mức: TH
Hai xe va chạm. Xe 1 chịu lực trung bình 7 $N$ từ xe 2. Khối lượng xe 1 là $3\,\text{kg}$. Tính độ lớn gia tốc trung bình của xe 1.
\choice
{$2,33 m/s^2$}
{\True $2,33 m/s^2$}
{$1,33 m/s^2$}
{$4,66 m/s^2$}
\loigiai{
Chọn B. Theo định luật II Newton, a=F/m=7/3= $2,33\,\text{m}$ /s^2; lực tương tác trên xe 2 có cùng độ lớn.
}
\end{ex}

% ===== Câu 89 =====
% ID: AIQ_L10C3_FULL_0575
% Mức: TH
\begin{ex}
% ID: AIQ_L10C3_FULL_0575
% Mức: TH
Hai xe va chạm. Xe 1 chịu lực trung bình 9 $N$ từ xe 2. Khối lượng xe 1 là $4\,\text{kg}$. Tính độ lớn gia tốc trung bình của xe 1.
\choice
{$2,25 m/s^2$}
{$3,25 m/s^2$}
{\True $2,25 m/s^2$}
{$4,5 m/s^2$}
\loigiai{
Chọn C. Theo định luật II Newton, a=F/m=9/4= $2,25\,\text{m}$ /s^2; lực tương tác trên xe 2 có cùng độ lớn.
}
\end{ex}

% ===== Câu 90 =====
% ID: AIQ_L10C3_FULL_0576
% Mức: VD
\begin{ex}
% ID: AIQ_L10C3_FULL_0576
% Mức: VD
Hai xe va chạm. Xe 1 chịu lực trung bình 11 $N$ từ xe 2. Khối lượng xe 1 là $5\,\text{kg}$. Tính độ lớn gia tốc trung bình của xe 1.
\choice
{$2,2 m/s^2$}
{$3,2 m/s^2$}
{$1,2 m/s^2$}
{\True $2,2 m/s^2$}
\loigiai{
Chọn D. Theo định luật II Newton, a=F/m=11/5= $2,2\,\text{m}$ /s^2; lực tương tác trên xe 2 có cùng độ lớn.
}
\end{ex}

% ===== Câu 91 =====
% ID: AIQ_L10C3_FULL_0577
% Mức: VD
\begin{ex}
% ID: AIQ_L10C3_FULL_0577
% Mức: VD
Hai xe va chạm. Xe 1 chịu lực trung bình 13 $N$ từ xe 2. Khối lượng xe 1 là $6\,\text{kg}$. Tính độ lớn gia tốc trung bình của xe 1.
\choice
{\True $2,17 m/s^2$}
{$3,17 m/s^2$}
{$1,17 m/s^2$}
{$4,34 m/s^2$}
\loigiai{
Chọn A. Theo định luật II Newton, a=F/m=13/6= $2,17\,\text{m}$ /s^2; lực tương tác trên xe 2 có cùng độ lớn.
}
\end{ex}

% ===== Câu 92 =====
% ID: AIQ_L10C3_FULL_0578
% Mức: TH
\begin{ex}
% ID: AIQ_L10C3_FULL_0578
% Mức: TH
Hai xe va chạm. Xe 1 chịu lực trung bình 5 $N$ từ xe 2. Khối lượng xe 1 là $5\,\text{kg}$. Tính độ lớn gia tốc trung bình của xe 1.
\choiceTF
{\True Giá trị cần tìm là $1\,\text{m}$ /s^2.}
{Giá trị cần tìm là $2\,\text{m}$ /s^2.}
{\True Phải xác định đúng các lực tác dụng và chọn chiều dương trước khi lập phương trình.}
{Có thể bỏ qua điều kiện áp dụng và chiều của lực mà kết quả vẫn luôn đúng.}
\loigiai{
\begin{itemchoice}
\item (Đúng) Giá trị cần tìm là $1\,\text{m}$ /s^2. (Vì): Theo định luật II Newton, a=F/m=5/5= $1\,\text{m}$ /s^2; lực tương tác trên xe 2 có cùng độ lớn. b) Sai vì sai giá trị. c) Đúng. d) Sai vì phải kiểm tra điều kiện áp dụng và chiều lực
\item (Sai) Giá trị cần tìm là $2\,\text{m}$ /s^2.
\item (Đúng) Phải xác định đúng các lực tác dụng và chọn chiều dương trước khi lập phương trình.
\item (Sai) Có thể bỏ qua điều kiện áp dụng và chiều của lực mà kết quả vẫn luôn đúng.
\end{itemchoice}
}
\end{ex}

% ===== Câu 93 =====
% ID: AIQ_L10C3_FULL_0579
% Mức: TH
\begin{ex}
% ID: AIQ_L10C3_FULL_0579
% Mức: TH
Hai xe va chạm. Xe 1 chịu lực trung bình 7 $N$ từ xe 2. Khối lượng xe 1 là $6\,\text{kg}$. Tính độ lớn gia tốc trung bình của xe 1.
\choiceTF
{\True Giá trị cần tìm là $1,17\,\text{m}$ /s^2.}
{Giá trị cần tìm là $2,17\,\text{m}$ /s^2.}
{\True Phải xác định đúng các lực tác dụng và chọn chiều dương trước khi lập phương trình.}
{Có thể bỏ qua điều kiện áp dụng và chiều của lực mà kết quả vẫn luôn đúng.}
\loigiai{
\begin{itemchoice}
\item (Đúng) Giá trị cần tìm là $1,17\,\text{m}$ /s^2. (Vì): Theo định luật II Newton, a=F/m=7/6= $1,17\,\text{m}$ /s^2; lực tương tác trên xe 2 có cùng độ lớn. b) Sai vì sai giá trị. c) Đúng. d) Sai vì phải kiểm tra điều kiện áp dụng và chiều lực
\item (Sai) Giá trị cần tìm là $2,17\,\text{m}$ /s^2.
\item (Đúng) Phải xác định đúng các lực tác dụng và chọn chiều dương trước khi lập phương trình.
\item (Sai) Có thể bỏ qua điều kiện áp dụng và chiều của lực mà kết quả vẫn luôn đúng.
\end{itemchoice}
}
\end{ex}

% ===== Câu 94 =====
% ID: AIQ_L10C3_FULL_0580
% Mức: VD
\begin{ex}
% ID: AIQ_L10C3_FULL_0580
% Mức: VD
Hai xe va chạm. Xe 1 chịu lực trung bình 9 $N$ từ xe 2. Khối lượng xe 1 là $2\,\text{kg}$. Tính độ lớn gia tốc trung bình của xe 1.
\choiceTF
{\True Giá trị cần tìm là $4,5\,\text{m}$ /s^2.}
{Giá trị cần tìm là $5,5\,\text{m}$ /s^2.}
{\True Phải xác định đúng các lực tác dụng và chọn chiều dương trước khi lập phương trình.}
{Có thể bỏ qua điều kiện áp dụng và chiều của lực mà kết quả vẫn luôn đúng.}
\loigiai{
\begin{itemchoice}
\item (Đúng) Giá trị cần tìm là $4,5\,\text{m}$ /s^2. (Vì): Theo định luật II Newton, a=F/m=9/2= $4,5\,\text{m}$ /s^2; lực tương tác trên xe 2 có cùng độ lớn. b) Sai vì sai giá trị. c) Đúng. d) Sai vì phải kiểm tra điều kiện áp dụng và chiều lực
\item (Sai) Giá trị cần tìm là $5,5\,\text{m}$ /s^2.
\item (Đúng) Phải xác định đúng các lực tác dụng và chọn chiều dương trước khi lập phương trình.
\item (Sai) Có thể bỏ qua điều kiện áp dụng và chiều của lực mà kết quả vẫn luôn đúng.
\end{itemchoice}
}
\end{ex}

% ===== Câu 95 =====
% ID: AIQ_L10C3_FULL_0581
% Mức: VD
\begin{ex}
% ID: AIQ_L10C3_FULL_0581
% Mức: VD
Hai xe va chạm. Xe 1 chịu lực trung bình 11 $N$ từ xe 2. Khối lượng xe 1 là $3\,\text{kg}$. Tính độ lớn gia tốc trung bình của xe 1.
\choiceTF
{\True Giá trị cần tìm là $3,67\,\text{m}$ /s^2.}
{Giá trị cần tìm là $4,67\,\text{m}$ /s^2.}
{\True Phải xác định đúng các lực tác dụng và chọn chiều dương trước khi lập phương trình.}
{Có thể bỏ qua điều kiện áp dụng và chiều của lực mà kết quả vẫn luôn đúng.}
\loigiai{
\begin{itemchoice}
\item (Đúng) Giá trị cần tìm là $3,67\,\text{m}$ /s^2. (Vì): Theo định luật II Newton, a=F/m=11/3= $3,67\,\text{m}$ /s^2; lực tương tác trên xe 2 có cùng độ lớn. b) Sai vì sai giá trị. c) Đúng. d) Sai vì phải kiểm tra điều kiện áp dụng và chiều lực
\item (Sai) Giá trị cần tìm là $4,67\,\text{m}$ /s^2.
\item (Đúng) Phải xác định đúng các lực tác dụng và chọn chiều dương trước khi lập phương trình.
\item (Sai) Có thể bỏ qua điều kiện áp dụng và chiều của lực mà kết quả vẫn luôn đúng.
\end{itemchoice}
}
\end{ex}

% ===== Câu 96 =====
% ID: AIQ_L10C3_FULL_0582
% Mức: VD
\begin{ex}
% ID: AIQ_L10C3_FULL_0582
% Mức: VD
Hai xe va chạm. Xe 1 chịu lực trung bình 13 $N$ từ xe 2. Khối lượng xe 1 là $4\,\text{kg}$. Tính độ lớn gia tốc trung bình của xe 1.
\choiceTF
{\True Giá trị cần tìm là $3,25\,\text{m}$ /s^2.}
{Giá trị cần tìm là $4,25\,\text{m}$ /s^2.}
{\True Phải xác định đúng các lực tác dụng và chọn chiều dương trước khi lập phương trình.}
{Có thể bỏ qua điều kiện áp dụng và chiều của lực mà kết quả vẫn luôn đúng.}
\loigiai{
\begin{itemchoice}
\item (Đúng) Giá trị cần tìm là $3,25\,\text{m}$ /s^2. (Vì): Theo định luật II Newton, a=F/m=13/4= $3,25\,\text{m}$ /s^2; lực tương tác trên xe 2 có cùng độ lớn. b) Sai vì sai giá trị. c) Đúng. d) Sai vì phải kiểm tra điều kiện áp dụng và chiều lực
\item (Sai) Giá trị cần tìm là $4,25\,\text{m}$ /s^2.
\item (Đúng) Phải xác định đúng các lực tác dụng và chọn chiều dương trước khi lập phương trình.
\item (Sai) Có thể bỏ qua điều kiện áp dụng và chiều của lực mà kết quả vẫn luôn đúng.
\end{itemchoice}
}
\end{ex}

% ===== Câu 97 =====
% ID: AIQ_L10C3_FULL_0583
% Mức: TH
\begin{ex}
% ID: AIQ_L10C3_FULL_0583
% Mức: TH
Hai xe va chạm. Xe 1 chịu lực trung bình 15 $N$ từ xe 2. Khối lượng xe 1 là $3\,\text{kg}$. Tính độ lớn gia tốc trung bình của xe 1.
\shortans{5}
\loigiai{
Theo định luật II Newton, a=F/m=15/3= $5\,\text{m}$ /s^2; lực tương tác trên xe 2 có cùng độ lớn.
}
\end{ex}

% ===== Câu 98 =====
% ID: AIQ_L10C3_FULL_0584
% Mức: TH
\begin{ex}
% ID: AIQ_L10C3_FULL_0584
% Mức: TH
Hai xe va chạm. Xe 1 chịu lực trung bình 17 $N$ từ xe 2. Khối lượng xe 1 là $4\,\text{kg}$. Tính độ lớn gia tốc trung bình của xe 1.
\shortans{4,25}
\loigiai{
Theo định luật II Newton, a=F/m=17/4= $4,25\,\text{m}$ /s^2; lực tương tác trên xe 2 có cùng độ lớn.
}
\end{ex}

% ===== Câu 99 =====
% ID: AIQ_L10C3_FULL_0585
% Mức: VD
\begin{ex}
% ID: AIQ_L10C3_FULL_0585
% Mức: VD
Hai xe va chạm. Xe 1 chịu lực trung bình 19 $N$ từ xe 2. Khối lượng xe 1 là $5\,\text{kg}$. Tính độ lớn gia tốc trung bình của xe 1.
\shortans{3,8}
\loigiai{
Theo định luật II Newton, a=F/m=19/5= $3,8\,\text{m}$ /s^2; lực tương tác trên xe 2 có cùng độ lớn.
}
\end{ex}

% ===== Câu 100 =====
% ID: AIQ_L10C3_FULL_0586
% Mức: VD
\begin{ex}
% ID: AIQ_L10C3_FULL_0586
% Mức: VD
Hai xe va chạm. Xe 1 chịu lực trung bình 5 $N$ từ xe 2. Khối lượng xe 1 là $6\,\text{kg}$. Tính độ lớn gia tốc trung bình của xe 1.
\shortans{0,83}
\loigiai{
Theo định luật II Newton, a=F/m=5/6= $0,83\,\text{m}$ /s^2; lực tương tác trên xe 2 có cùng độ lớn.
}
\end{ex}

% ===== Câu 101 =====
% ID: AIQ_L10C3_FULL_0587
% Mức: VD
\begin{ex}
% ID: AIQ_L10C3_FULL_0587
% Mức: VD
Hai xe va chạm. Xe 1 chịu lực trung bình 7 $N$ từ xe 2. Khối lượng xe 1 là $2\,\text{kg}$. Tính độ lớn gia tốc trung bình của xe 1.
\shortans{3,5}
\loigiai{
Theo định luật II Newton, a=F/m=7/2= $3,5\,\text{m}$ /s^2; lực tương tác trên xe 2 có cùng độ lớn.
}
\end{ex}

% ===== Câu 102 =====
% ID: AIQ_L10C3_FULL_0588
% Mức: VDC
\begin{ex}
% ID: AIQ_L10C3_FULL_0588
% Mức: VDC
Hai xe va chạm. Xe 1 chịu lực trung bình 9 $N$ từ xe 2. Khối lượng xe 1 là $3\,\text{kg}$. Tính độ lớn gia tốc trung bình của xe 1.
\shortans{3}
\loigiai{
Theo định luật II Newton, a=F/m=9/3= $3\,\text{m}$ /s^2; lực tương tác trên xe 2 có cùng độ lớn.
}
\end{ex}

% ===== Câu 103 =====
% ID: AIQ_L10C3_FULL_0589
% Mức: VD
\begin{bt}
% ID: AIQ_L10C3_FULL_0589
% Mức: VD
Hai xe va chạm. Xe 1 chịu lực trung bình 9 $N$ từ xe 2. Khối lượng xe 1 là $6\,\text{kg}$. Tính độ lớn gia tốc trung bình của xe 1.
\loigiai{
Phân tích lực, chọn hệ quy chiếu và áp dụng định luật hoặc quy tắc phù hợp. Theo định luật II Newton, a=F/m=9/6= $1,5\,\text{m}$ /s^2; lực tương tác trên xe 2 có cùng độ lớn.
}
\end{bt}

% ===== Câu 104 =====
% ID: AIQ_L10C3_FULL_0590
% Mức: VD
\begin{bt}
% ID: AIQ_L10C3_FULL_0590
% Mức: VD
Hai xe va chạm. Xe 1 chịu lực trung bình 11 $N$ từ xe 2. Khối lượng xe 1 là $2\,\text{kg}$. Tính độ lớn gia tốc trung bình của xe 1.
\loigiai{
Phân tích lực, chọn hệ quy chiếu và áp dụng định luật hoặc quy tắc phù hợp. Theo định luật II Newton, a=F/m=11/2= $5,5\,\text{m}$ /s^2; lực tương tác trên xe 2 có cùng độ lớn.
}
\end{bt}

% ===== Câu 105 =====
% ID: AIQ_L10C3_FULL_0591
% Mức: VDC
\begin{bt}
% ID: AIQ_L10C3_FULL_0591
% Mức: VDC
Hai xe va chạm. Xe 1 chịu lực trung bình 13 $N$ từ xe 2. Khối lượng xe 1 là $3\,\text{kg}$. Tính độ lớn gia tốc trung bình của xe 1.
\loigiai{
Phân tích lực, chọn hệ quy chiếu và áp dụng định luật hoặc quy tắc phù hợp. Theo định luật II Newton, a=F/m=13/3= $4,33\,\text{m}$ /s^2; lực tương tác trên xe 2 có cùng độ lớn.
}
\end{bt}

% ===== Câu 106 =====
% ID: AIQ_L10C3_FULL_0592
% Mức: VDC
\begin{bt}
% ID: AIQ_L10C3_FULL_0592
% Mức: VDC
Hai xe va chạm. Xe 1 chịu lực trung bình 15 $N$ từ xe 2. Khối lượng xe 1 là $4\,\text{kg}$. Tính độ lớn gia tốc trung bình của xe 1.
\loigiai{
Phân tích lực, chọn hệ quy chiếu và áp dụng định luật hoặc quy tắc phù hợp. Theo định luật II Newton, a=F/m=15/4= $3,75\,\text{m}$ /s^2; lực tương tác trên xe 2 có cùng độ lớn.
}
\end{bt}

% ===== Câu 107 =====
% ID: AIQ_L10C3_FULL_0593
% Mức: VDC
\begin{bt}
% ID: AIQ_L10C3_FULL_0593
% Mức: VDC
Hai xe va chạm. Xe 1 chịu lực trung bình 17 $N$ từ xe 2. Khối lượng xe 1 là $5\,\text{kg}$. Tính độ lớn gia tốc trung bình của xe 1.
\loigiai{
Phân tích lực, chọn hệ quy chiếu và áp dụng định luật hoặc quy tắc phù hợp. Theo định luật II Newton, a=F/m=17/5= $3,4\,\text{m}$ /s^2; lực tương tác trên xe 2 có cùng độ lớn.
}
\end{bt}

% ===== Câu 108 =====
% ID: AIQ_L10C3_FULL_0594
% Mức: VDC
\begin{bt}
% ID: AIQ_L10C3_FULL_0594
% Mức: VDC
Hai xe va chạm. Xe 1 chịu lực trung bình 19 $N$ từ xe 2. Khối lượng xe 1 là $6\,\text{kg}$. Tính độ lớn gia tốc trung bình của xe 1.
\loigiai{
Phân tích lực, chọn hệ quy chiếu và áp dụng định luật hoặc quy tắc phù hợp. Theo định luật II Newton, a=F/m=19/6= $3,17\,\text{m}$ /s^2; lực tương tác trên xe 2 có cùng độ lớn.
}
\end{bt}

% ===== Câu 109 =====
% ID: AIQ_L10C3_FULL_0595
% Mức: NB
\dangbt{Giải thích chuyển động bằng lực và phản lực}
\begin{ex}
% ID: AIQ_L10C3_FULL_0595
% Mức: NB
Người bơi đẩy nước về phía sau một lực $13\,\text{N}$. Bỏ qua lực cản khác, nước tác dụng lên người lực đẩy về phía trước bằng bao nhiêu?
\choice
{\True 13 $N$}
{14 $N$}
{12 $N$}
{26 $N$}
\loigiai{
Chọn A. Theo định luật III Newton, nước đẩy người về phía trước một lực bằng $13\,\text{N}$.
}
\end{ex}

% ===== Câu 110 =====
% ID: AIQ_L10C3_FULL_0596
% Mức: NB
\begin{ex}
% ID: AIQ_L10C3_FULL_0596
% Mức: NB
Người bơi đẩy nước về phía sau một lực $15\,\text{N}$. Bỏ qua lực cản khác, nước tác dụng lên người lực đẩy về phía trước bằng bao nhiêu?
\choice
{15 $N$}
{\True 15 $N$}
{14 $N$}
{30 $N$}
\loigiai{
Chọn B. Theo định luật III Newton, nước đẩy người về phía trước một lực bằng $15\,\text{N}$.
}
\end{ex}

% ===== Câu 111 =====
% ID: AIQ_L10C3_FULL_0597
% Mức: NB
\begin{ex}
% ID: AIQ_L10C3_FULL_0597
% Mức: NB
Người bơi đẩy nước về phía sau một lực $17\,\text{N}$. Bỏ qua lực cản khác, nước tác dụng lên người lực đẩy về phía trước bằng bao nhiêu?
\choice
{17 $N$}
{18 $N$}
{\True 17 $N$}
{34 $N$}
\loigiai{
Chọn C. Theo định luật III Newton, nước đẩy người về phía trước một lực bằng $17\,\text{N}$.
}
\end{ex}

% ===== Câu 112 =====
% ID: AIQ_L10C3_FULL_0598
% Mức: NB
\begin{ex}
% ID: AIQ_L10C3_FULL_0598
% Mức: NB
Người bơi đẩy nước về phía sau một lực $19\,\text{N}$. Bỏ qua lực cản khác, nước tác dụng lên người lực đẩy về phía trước bằng bao nhiêu?
\choice
{19 $N$}
{20 $N$}
{18 $N$}
{\True 19 $N$}
\loigiai{
Chọn D. Theo định luật III Newton, nước đẩy người về phía trước một lực bằng $19\,\text{N}$.
}
\end{ex}

% ===== Câu 113 =====
% ID: AIQ_L10C3_FULL_0599
% Mức: TH
\begin{ex}
% ID: AIQ_L10C3_FULL_0599
% Mức: TH
Người bơi đẩy nước về phía sau một lực $5\,\text{N}$. Bỏ qua lực cản khác, nước tác dụng lên người lực đẩy về phía trước bằng bao nhiêu?
\choice
{\True 5 $N$}
{6 $N$}
{4 $N$}
{10 $N$}
\loigiai{
Chọn A. Theo định luật III Newton, nước đẩy người về phía trước một lực bằng $5\,\text{N}$.
}
\end{ex}

% ===== Câu 114 =====
% ID: AIQ_L10C3_FULL_0600
% Mức: TH
\begin{ex}
% ID: AIQ_L10C3_FULL_0600
% Mức: TH
Người bơi đẩy nước về phía sau một lực $7\,\text{N}$. Bỏ qua lực cản khác, nước tác dụng lên người lực đẩy về phía trước bằng bao nhiêu?
\choice
{7 $N$}
{\True 7 $N$}
{6 $N$}
{14 $N$}
\loigiai{
Chọn B. Theo định luật III Newton, nước đẩy người về phía trước một lực bằng $7\,\text{N}$.
}
\end{ex}

% ===== Câu 115 =====
% ID: AIQ_L10C3_FULL_0601
% Mức: TH
\begin{ex}
% ID: AIQ_L10C3_FULL_0601
% Mức: TH
Người bơi đẩy nước về phía sau một lực $9\,\text{N}$. Bỏ qua lực cản khác, nước tác dụng lên người lực đẩy về phía trước bằng bao nhiêu?
\choice
{9 $N$}
{10 $N$}
{\True 9 $N$}
{18 $N$}
\loigiai{
Chọn C. Theo định luật III Newton, nước đẩy người về phía trước một lực bằng $9\,\text{N}$.
}
\end{ex}

% ===== Câu 116 =====
% ID: AIQ_L10C3_FULL_0602
% Mức: TH
\begin{ex}
% ID: AIQ_L10C3_FULL_0602
% Mức: TH
Người bơi đẩy nước về phía sau một lực $11\,\text{N}$. Bỏ qua lực cản khác, nước tác dụng lên người lực đẩy về phía trước bằng bao nhiêu?
\choice
{11 $N$}
{12 $N$}
{10 $N$}
{\True 11 $N$}
\loigiai{
Chọn D. Theo định luật III Newton, nước đẩy người về phía trước một lực bằng $11\,\text{N}$.
}
\end{ex}

% ===== Câu 117 =====
% ID: AIQ_L10C3_FULL_0603
% Mức: VD
\begin{ex}
% ID: AIQ_L10C3_FULL_0603
% Mức: VD
Người bơi đẩy nước về phía sau một lực $13\,\text{N}$. Bỏ qua lực cản khác, nước tác dụng lên người lực đẩy về phía trước bằng bao nhiêu?
\choice
{\True 13 $N$}
{14 $N$}
{12 $N$}
{26 $N$}
\loigiai{
Chọn A. Theo định luật III Newton, nước đẩy người về phía trước một lực bằng $13\,\text{N}$.
}
\end{ex}

% ===== Câu 118 =====
% ID: AIQ_L10C3_FULL_0604
% Mức: VD
\begin{ex}
% ID: AIQ_L10C3_FULL_0604
% Mức: VD
Người bơi đẩy nước về phía sau một lực $15\,\text{N}$. Bỏ qua lực cản khác, nước tác dụng lên người lực đẩy về phía trước bằng bao nhiêu?
\choice
{15 $N$}
{\True 15 $N$}
{14 $N$}
{30 $N$}
\loigiai{
Chọn B. Theo định luật III Newton, nước đẩy người về phía trước một lực bằng $15\,\text{N}$.
}
\end{ex}

% ===== Câu 119 =====
% ID: AIQ_L10C3_FULL_0605
% Mức: TH
\begin{ex}
% ID: AIQ_L10C3_FULL_0605
% Mức: TH
Người bơi đẩy nước về phía sau một lực $7\,\text{N}$. Bỏ qua lực cản khác, nước tác dụng lên người lực đẩy về phía trước bằng bao nhiêu?
\choiceTF
{\True Giá trị cần tìm là $7\,\text{N}$.}
{Giá trị cần tìm là $8\,\text{N}$.}
{\True Phải xác định đúng các lực tác dụng và chọn chiều dương trước khi lập phương trình.}
{Có thể bỏ qua điều kiện áp dụng và chiều của lực mà kết quả vẫn luôn đúng.}
\loigiai{
\begin{itemchoice}
\item (Đúng) Giá trị cần tìm là $7\,\text{N}$. (Vì): Theo định luật III Newton, nước đẩy người về phía trước một lực bằng $7\,\text{N}$. b) Sai vì sai giá trị. c) Đúng. d) Sai vì phải kiểm tra điều kiện áp dụng và chiều lực
\item (Sai) Giá trị cần tìm là $8\,\text{N}$.
\item (Đúng) Phải xác định đúng các lực tác dụng và chọn chiều dương trước khi lập phương trình.
\item (Sai) Có thể bỏ qua điều kiện áp dụng và chiều của lực mà kết quả vẫn luôn đúng.
\end{itemchoice}
}
\end{ex}

% ===== Câu 120 =====
% ID: AIQ_L10C3_FULL_0606
% Mức: TH
\begin{ex}
% ID: AIQ_L10C3_FULL_0606
% Mức: TH
Người bơi đẩy nước về phía sau một lực $9\,\text{N}$. Bỏ qua lực cản khác, nước tác dụng lên người lực đẩy về phía trước bằng bao nhiêu?
\choiceTF
{\True Giá trị cần tìm là $9\,\text{N}$.}
{Giá trị cần tìm là $10\,\text{N}$.}
{\True Phải xác định đúng các lực tác dụng và chọn chiều dương trước khi lập phương trình.}
{Có thể bỏ qua điều kiện áp dụng và chiều của lực mà kết quả vẫn luôn đúng.}
\loigiai{
\begin{itemchoice}
\item (Đúng) Giá trị cần tìm là $9\,\text{N}$. (Vì): Theo định luật III Newton, nước đẩy người về phía trước một lực bằng $9\,\text{N}$. b) Sai vì sai giá trị. c) Đúng. d) Sai vì phải kiểm tra điều kiện áp dụng và chiều lực
\item (Sai) Giá trị cần tìm là $10\,\text{N}$.
\item (Đúng) Phải xác định đúng các lực tác dụng và chọn chiều dương trước khi lập phương trình.
\item (Sai) Có thể bỏ qua điều kiện áp dụng và chiều của lực mà kết quả vẫn luôn đúng.
\end{itemchoice}
}
\end{ex}

% ===== Câu 121 =====
% ID: AIQ_L10C3_FULL_0607
% Mức: VD
\begin{ex}
% ID: AIQ_L10C3_FULL_0607
% Mức: VD
Người bơi đẩy nước về phía sau một lực $11\,\text{N}$. Bỏ qua lực cản khác, nước tác dụng lên người lực đẩy về phía trước bằng bao nhiêu?
\choiceTF
{\True Giá trị cần tìm là $11\,\text{N}$.}
{Giá trị cần tìm là $12\,\text{N}$.}
{\True Phải xác định đúng các lực tác dụng và chọn chiều dương trước khi lập phương trình.}
{Có thể bỏ qua điều kiện áp dụng và chiều của lực mà kết quả vẫn luôn đúng.}
\loigiai{
\begin{itemchoice}
\item (Đúng) Giá trị cần tìm là $11\,\text{N}$. (Vì): Theo định luật III Newton, nước đẩy người về phía trước một lực bằng $11\,\text{N}$. b) Sai vì sai giá trị. c) Đúng. d) Sai vì phải kiểm tra điều kiện áp dụng và chiều lực
\item (Sai) Giá trị cần tìm là $12\,\text{N}$.
\item (Đúng) Phải xác định đúng các lực tác dụng và chọn chiều dương trước khi lập phương trình.
\item (Sai) Có thể bỏ qua điều kiện áp dụng và chiều của lực mà kết quả vẫn luôn đúng.
\end{itemchoice}
}
\end{ex}

% ===== Câu 122 =====
% ID: AIQ_L10C3_FULL_0608
% Mức: VD
\begin{ex}
% ID: AIQ_L10C3_FULL_0608
% Mức: VD
Người bơi đẩy nước về phía sau một lực $13\,\text{N}$. Bỏ qua lực cản khác, nước tác dụng lên người lực đẩy về phía trước bằng bao nhiêu?
\choiceTF
{\True Giá trị cần tìm là $13\,\text{N}$.}
{Giá trị cần tìm là $14\,\text{N}$.}
{\True Phải xác định đúng các lực tác dụng và chọn chiều dương trước khi lập phương trình.}
{Có thể bỏ qua điều kiện áp dụng và chiều của lực mà kết quả vẫn luôn đúng.}
\loigiai{
\begin{itemchoice}
\item (Đúng) Giá trị cần tìm là $13\,\text{N}$. (Vì): Theo định luật III Newton, nước đẩy người về phía trước một lực bằng $13\,\text{N}$. b) Sai vì sai giá trị. c) Đúng. d) Sai vì phải kiểm tra điều kiện áp dụng và chiều lực
\item (Sai) Giá trị cần tìm là $14\,\text{N}$.
\item (Đúng) Phải xác định đúng các lực tác dụng và chọn chiều dương trước khi lập phương trình.
\item (Sai) Có thể bỏ qua điều kiện áp dụng và chiều của lực mà kết quả vẫn luôn đúng.
\end{itemchoice}
}
\end{ex}

% ===== Câu 123 =====
% ID: AIQ_L10C3_FULL_0609
% Mức: VD
\begin{ex}
% ID: AIQ_L10C3_FULL_0609
% Mức: VD
Người bơi đẩy nước về phía sau một lực $15\,\text{N}$. Bỏ qua lực cản khác, nước tác dụng lên người lực đẩy về phía trước bằng bao nhiêu?
\choiceTF
{\True Giá trị cần tìm là $15\,\text{N}$.}
{Giá trị cần tìm là $16\,\text{N}$.}
{\True Phải xác định đúng các lực tác dụng và chọn chiều dương trước khi lập phương trình.}
{Có thể bỏ qua điều kiện áp dụng và chiều của lực mà kết quả vẫn luôn đúng.}
\loigiai{
\begin{itemchoice}
\item (Đúng) Giá trị cần tìm là $15\,\text{N}$. (Vì): Theo định luật III Newton, nước đẩy người về phía trước một lực bằng $15\,\text{N}$. b) Sai vì sai giá trị. c) Đúng. d) Sai vì phải kiểm tra điều kiện áp dụng và chiều lực
\item (Sai) Giá trị cần tìm là $16\,\text{N}$.
\item (Đúng) Phải xác định đúng các lực tác dụng và chọn chiều dương trước khi lập phương trình.
\item (Sai) Có thể bỏ qua điều kiện áp dụng và chiều của lực mà kết quả vẫn luôn đúng.
\end{itemchoice}
}
\end{ex}

% ===== Câu 124 =====
% ID: AIQ_L10C3_FULL_0610
% Mức: TH
\begin{ex}
% ID: AIQ_L10C3_FULL_0610
% Mức: TH
Người bơi đẩy nước về phía sau một lực $17\,\text{N}$. Bỏ qua lực cản khác, nước tác dụng lên người lực đẩy về phía trước bằng bao nhiêu?
\shortans{17}
\loigiai{
Theo định luật III Newton, nước đẩy người về phía trước một lực bằng $17\,\text{N}$.
}
\end{ex}

% ===== Câu 125 =====
% ID: AIQ_L10C3_FULL_0611
% Mức: TH
\begin{ex}
% ID: AIQ_L10C3_FULL_0611
% Mức: TH
Người bơi đẩy nước về phía sau một lực $19\,\text{N}$. Bỏ qua lực cản khác, nước tác dụng lên người lực đẩy về phía trước bằng bao nhiêu?
\shortans{19}
\loigiai{
Theo định luật III Newton, nước đẩy người về phía trước một lực bằng $19\,\text{N}$.
}
\end{ex}

% ===== Câu 126 =====
% ID: AIQ_L10C3_FULL_0612
% Mức: VD
\begin{ex}
% ID: AIQ_L10C3_FULL_0612
% Mức: VD
Người bơi đẩy nước về phía sau một lực $5\,\text{N}$. Bỏ qua lực cản khác, nước tác dụng lên người lực đẩy về phía trước bằng bao nhiêu?
\shortans{5}
\loigiai{
Theo định luật III Newton, nước đẩy người về phía trước một lực bằng $5\,\text{N}$.
}
\end{ex}

% ===== Câu 127 =====
% ID: AIQ_L10C3_FULL_0613
% Mức: VD
\begin{ex}
% ID: AIQ_L10C3_FULL_0613
% Mức: VD
Người bơi đẩy nước về phía sau một lực $7\,\text{N}$. Bỏ qua lực cản khác, nước tác dụng lên người lực đẩy về phía trước bằng bao nhiêu?
\shortans{7}
\loigiai{
Theo định luật III Newton, nước đẩy người về phía trước một lực bằng $7\,\text{N}$.
}
\end{ex}

% ===== Câu 128 =====
% ID: AIQ_L10C3_FULL_0614
% Mức: VD
\begin{ex}
% ID: AIQ_L10C3_FULL_0614
% Mức: VD
Người bơi đẩy nước về phía sau một lực $9\,\text{N}$. Bỏ qua lực cản khác, nước tác dụng lên người lực đẩy về phía trước bằng bao nhiêu?
\shortans{9}
\loigiai{
Theo định luật III Newton, nước đẩy người về phía trước một lực bằng $9\,\text{N}$.
}
\end{ex}

% ===== Câu 129 =====
% ID: AIQ_L10C3_FULL_0615
% Mức: VDC
\begin{ex}
% ID: AIQ_L10C3_FULL_0615
% Mức: VDC
Người bơi đẩy nước về phía sau một lực $11\,\text{N}$. Bỏ qua lực cản khác, nước tác dụng lên người lực đẩy về phía trước bằng bao nhiêu?
\shortans{11}
\loigiai{
Theo định luật III Newton, nước đẩy người về phía trước một lực bằng $11\,\text{N}$.
}
\end{ex}

% ===== Câu 130 =====
% ID: AIQ_L10C3_FULL_0616
% Mức: VD
\begin{bt}
% ID: AIQ_L10C3_FULL_0616
% Mức: VD
Người bơi đẩy nước về phía sau một lực $11\,\text{N}$. Bỏ qua lực cản khác, nước tác dụng lên người lực đẩy về phía trước bằng bao nhiêu?
\loigiai{
Phân tích lực, chọn hệ quy chiếu và áp dụng định luật hoặc quy tắc phù hợp. Theo định luật III Newton, nước đẩy người về phía trước một lực bằng $11\,\text{N}$.
}
\end{bt}

% ===== Câu 131 =====
% ID: AIQ_L10C3_FULL_0617
% Mức: VD
\begin{bt}
% ID: AIQ_L10C3_FULL_0617
% Mức: VD
Người bơi đẩy nước về phía sau một lực $13\,\text{N}$. Bỏ qua lực cản khác, nước tác dụng lên người lực đẩy về phía trước bằng bao nhiêu?
\loigiai{
Phân tích lực, chọn hệ quy chiếu và áp dụng định luật hoặc quy tắc phù hợp. Theo định luật III Newton, nước đẩy người về phía trước một lực bằng $13\,\text{N}$.
}
\end{bt}

% ===== Câu 132 =====
% ID: AIQ_L10C3_FULL_0618
% Mức: VDC
\begin{bt}
% ID: AIQ_L10C3_FULL_0618
% Mức: VDC
Người bơi đẩy nước về phía sau một lực $15\,\text{N}$. Bỏ qua lực cản khác, nước tác dụng lên người lực đẩy về phía trước bằng bao nhiêu?
\loigiai{
Phân tích lực, chọn hệ quy chiếu và áp dụng định luật hoặc quy tắc phù hợp. Theo định luật III Newton, nước đẩy người về phía trước một lực bằng $15\,\text{N}$.
}
\end{bt}

% ===== Câu 133 =====
% ID: AIQ_L10C3_FULL_0619
% Mức: VDC
\begin{bt}
% ID: AIQ_L10C3_FULL_0619
% Mức: VDC
Người bơi đẩy nước về phía sau một lực $17\,\text{N}$. Bỏ qua lực cản khác, nước tác dụng lên người lực đẩy về phía trước bằng bao nhiêu?
\loigiai{
Phân tích lực, chọn hệ quy chiếu và áp dụng định luật hoặc quy tắc phù hợp. Theo định luật III Newton, nước đẩy người về phía trước một lực bằng $17\,\text{N}$.
}
\end{bt}

% ===== Câu 134 =====
% ID: AIQ_L10C3_FULL_0620
% Mức: VDC
\begin{bt}
% ID: AIQ_L10C3_FULL_0620
% Mức: VDC
Người bơi đẩy nước về phía sau một lực $19\,\text{N}$. Bỏ qua lực cản khác, nước tác dụng lên người lực đẩy về phía trước bằng bao nhiêu?
\loigiai{
Phân tích lực, chọn hệ quy chiếu và áp dụng định luật hoặc quy tắc phù hợp. Theo định luật III Newton, nước đẩy người về phía trước một lực bằng $19\,\text{N}$.
}
\end{bt}

% ===== Câu 135 =====
% ID: AIQ_L10C3_FULL_0621
% Mức: VDC
\begin{bt}
% ID: AIQ_L10C3_FULL_0621
% Mức: VDC
Người bơi đẩy nước về phía sau một lực $5\,\text{N}$. Bỏ qua lực cản khác, nước tác dụng lên người lực đẩy về phía trước bằng bao nhiêu?
\loigiai{
Phân tích lực, chọn hệ quy chiếu và áp dụng định luật hoặc quy tắc phù hợp. Theo định luật III Newton, nước đẩy người về phía trước một lực bằng $5\,\text{N}$.
}
\end{bt}

% ===== Câu 136 =====
% ID: AIQ_L10C3_FULL_0622
% Mức: NB
\dangbt{Bài toán tổng hợp nhiều vật tương tác}
\begin{ex}
% ID: AIQ_L10C3_FULL_0622
% Mức: NB
Hai vật tiếp xúc cùng chuyển động, hợp lực ngoài làm hệ có gia tốc $2\,\text{m}$ /s^2. Tổng khối lượng hệ là $10\,\text{kg}$. Tính hợp lực ngoài tác dụng lên hệ.
\choice
{20 $N$}
{\True 20 $N$}
{19 $N$}
{40 $N$}
\loigiai{
Chọn B. Xét toàn hệ, các lực tương tác trong triệt tiêu theo từng cặp; F=(m1+m2)a=10·2= $20\,\text{N}$.
}
\end{ex}

% ===== Câu 137 =====
% ID: AIQ_L10C3_FULL_0623
% Mức: NB
\begin{ex}
% ID: AIQ_L10C3_FULL_0623
% Mức: NB
Hai vật tiếp xúc cùng chuyển động, hợp lực ngoài làm hệ có gia tốc $3\,\text{m}$ /s^2. Tổng khối lượng hệ là $7\,\text{kg}$. Tính hợp lực ngoài tác dụng lên hệ.
\choice
{21 $N$}
{22 $N$}
{\True 21 $N$}
{42 $N$}
\loigiai{
Chọn C. Xét toàn hệ, các lực tương tác trong triệt tiêu theo từng cặp; F=(m1+m2)a=7·3= $21\,\text{N}$.
}
\end{ex}

% ===== Câu 138 =====
% ID: AIQ_L10C3_FULL_0624
% Mức: NB
\begin{ex}
% ID: AIQ_L10C3_FULL_0624
% Mức: NB
Hai vật tiếp xúc cùng chuyển động, hợp lực ngoài làm hệ có gia tốc $4\,\text{m}$ /s^2. Tổng khối lượng hệ là $10\,\text{kg}$. Tính hợp lực ngoài tác dụng lên hệ.
\choice
{40 $N$}
{41 $N$}
{39 $N$}
{\True 40 $N$}
\loigiai{
Chọn D. Xét toàn hệ, các lực tương tác trong triệt tiêu theo từng cặp; F=(m1+m2)a=10·4= $40\,\text{N}$.
}
\end{ex}

% ===== Câu 139 =====
% ID: AIQ_L10C3_FULL_0625
% Mức: NB
\begin{ex}
% ID: AIQ_L10C3_FULL_0625
% Mức: NB
Hai vật tiếp xúc cùng chuyển động, hợp lực ngoài làm hệ có gia tốc $1\,\text{m}$ /s^2. Tổng khối lượng hệ là $13\,\text{kg}$. Tính hợp lực ngoài tác dụng lên hệ.
\choice
{\True 13 $N$}
{14 $N$}
{12 $N$}
{26 $N$}
\loigiai{
Chọn A. Xét toàn hệ, các lực tương tác trong triệt tiêu theo từng cặp; F=(m1+m2)a=13·1= $13\,\text{N}$.
}
\end{ex}

% ===== Câu 140 =====
% ID: AIQ_L10C3_FULL_0626
% Mức: TH
\begin{ex}
% ID: AIQ_L10C3_FULL_0626
% Mức: TH
Hai vật tiếp xúc cùng chuyển động, hợp lực ngoài làm hệ có gia tốc $2\,\text{m}$ /s^2. Tổng khối lượng hệ là $10\,\text{kg}$. Tính hợp lực ngoài tác dụng lên hệ.
\choice
{20 $N$}
{\True 20 $N$}
{19 $N$}
{40 $N$}
\loigiai{
Chọn B. Xét toàn hệ, các lực tương tác trong triệt tiêu theo từng cặp; F=(m1+m2)a=10·2= $20\,\text{N}$.
}
\end{ex}

% ===== Câu 141 =====
% ID: AIQ_L10C3_FULL_0627
% Mức: TH
\begin{ex}
% ID: AIQ_L10C3_FULL_0627
% Mức: TH
Hai vật tiếp xúc cùng chuyển động, hợp lực ngoài làm hệ có gia tốc $3\,\text{m}$ /s^2. Tổng khối lượng hệ là $8\,\text{kg}$. Tính hợp lực ngoài tác dụng lên hệ.
\choice
{24 $N$}
{25 $N$}
{\True 24 $N$}
{48 $N$}
\loigiai{
Chọn C. Xét toàn hệ, các lực tương tác trong triệt tiêu theo từng cặp; F=(m1+m2)a=8·3= $24\,\text{N}$.
}
\end{ex}

% ===== Câu 142 =====
% ID: AIQ_L10C3_FULL_0628
% Mức: TH
\begin{ex}
% ID: AIQ_L10C3_FULL_0628
% Mức: TH
Hai vật tiếp xúc cùng chuyển động, hợp lực ngoài làm hệ có gia tốc $4\,\text{m}$ /s^2. Tổng khối lượng hệ là $11\,\text{kg}$. Tính hợp lực ngoài tác dụng lên hệ.
\choice
{44 $N$}
{45 $N$}
{43 $N$}
{\True 44 $N$}
\loigiai{
Chọn D. Xét toàn hệ, các lực tương tác trong triệt tiêu theo từng cặp; F=(m1+m2)a=11·4= $44\,\text{N}$.
}
\end{ex}

% ===== Câu 143 =====
% ID: AIQ_L10C3_FULL_0629
% Mức: TH
\begin{ex}
% ID: AIQ_L10C3_FULL_0629
% Mức: TH
Hai vật tiếp xúc cùng chuyển động, hợp lực ngoài làm hệ có gia tốc $1\,\text{m}$ /s^2. Tổng khối lượng hệ là $8\,\text{kg}$. Tính hợp lực ngoài tác dụng lên hệ.
\choice
{\True 8 $N$}
{9 $N$}
{7 $N$}
{16 $N$}
\loigiai{
Chọn A. Xét toàn hệ, các lực tương tác trong triệt tiêu theo từng cặp; F=(m1+m2)a=8·1= $8\,\text{N}$.
}
\end{ex}

% ===== Câu 144 =====
% ID: AIQ_L10C3_FULL_0630
% Mức: VD
\begin{ex}
% ID: AIQ_L10C3_FULL_0630
% Mức: VD
Hai vật tiếp xúc cùng chuyển động, hợp lực ngoài làm hệ có gia tốc $2\,\text{m}$ /s^2. Tổng khối lượng hệ là $11\,\text{kg}$. Tính hợp lực ngoài tác dụng lên hệ.
\choice
{22 $N$}
{\True 22 $N$}
{21 $N$}
{44 $N$}
\loigiai{
Chọn B. Xét toàn hệ, các lực tương tác trong triệt tiêu theo từng cặp; F=(m1+m2)a=11·2= $22\,\text{N}$.
}
\end{ex}

% ===== Câu 145 =====
% ID: AIQ_L10C3_FULL_0631
% Mức: VD
\begin{ex}
% ID: AIQ_L10C3_FULL_0631
% Mức: VD
Hai vật tiếp xúc cùng chuyển động, hợp lực ngoài làm hệ có gia tốc $3\,\text{m}$ /s^2. Tổng khối lượng hệ là $14\,\text{kg}$. Tính hợp lực ngoài tác dụng lên hệ.
\choice
{42 $N$}
{43 $N$}
{\True 42 $N$}
{84 $N$}
\loigiai{
Chọn C. Xét toàn hệ, các lực tương tác trong triệt tiêu theo từng cặp; F=(m1+m2)a=14·3= $42\,\text{N}$.
}
\end{ex}

% ===== Câu 146 =====
% ID: AIQ_L10C3_FULL_0632
% Mức: TH
\begin{ex}
% ID: AIQ_L10C3_FULL_0632
% Mức: TH
Hai vật tiếp xúc cùng chuyển động, hợp lực ngoài làm hệ có gia tốc $3\,\text{m}$ /s^2. Tổng khối lượng hệ là $9\,\text{kg}$. Tính hợp lực ngoài tác dụng lên hệ.
\choiceTF
{\True Giá trị cần tìm là $27\,\text{N}$.}
{Giá trị cần tìm là $28\,\text{N}$.}
{\True Phải xác định đúng các lực tác dụng và chọn chiều dương trước khi lập phương trình.}
{Có thể bỏ qua điều kiện áp dụng và chiều của lực mà kết quả vẫn luôn đúng.}
\loigiai{
\begin{itemchoice}
\item (Đúng) Giá trị cần tìm là $27\,\text{N}$. (Vì): Xét toàn hệ, các lực tương tác trong triệt tiêu theo từng cặp; F=(m1+m2)a=9·3= $27\,\text{N}$. b) Sai vì sai giá trị. c) Đúng. d) Sai vì phải kiểm tra điều kiện áp dụng và chiều lực
\item (Sai) Giá trị cần tìm là $28\,\text{N}$.
\item (Đúng) Phải xác định đúng các lực tác dụng và chọn chiều dương trước khi lập phương trình.
\item (Sai) Có thể bỏ qua điều kiện áp dụng và chiều của lực mà kết quả vẫn luôn đúng.
\end{itemchoice}
}
\end{ex}

% ===== Câu 147 =====
% ID: AIQ_L10C3_FULL_0633
% Mức: TH
\begin{ex}
% ID: AIQ_L10C3_FULL_0633
% Mức: TH
Hai vật tiếp xúc cùng chuyển động, hợp lực ngoài làm hệ có gia tốc $4\,\text{m}$ /s^2. Tổng khối lượng hệ là $12\,\text{kg}$. Tính hợp lực ngoài tác dụng lên hệ.
\choiceTF
{\True Giá trị cần tìm là $48\,\text{N}$.}
{Giá trị cần tìm là $49\,\text{N}$.}
{\True Phải xác định đúng các lực tác dụng và chọn chiều dương trước khi lập phương trình.}
{Có thể bỏ qua điều kiện áp dụng và chiều của lực mà kết quả vẫn luôn đúng.}
\loigiai{
\begin{itemchoice}
\item (Đúng) Giá trị cần tìm là $48\,\text{N}$. (Vì): Xét toàn hệ, các lực tương tác trong triệt tiêu theo từng cặp; F=(m1+m2)a=12·4= $48\,\text{N}$. b) Sai vì sai giá trị. c) Đúng. d) Sai vì phải kiểm tra điều kiện áp dụng và chiều lực
\item (Sai) Giá trị cần tìm là $49\,\text{N}$.
\item (Đúng) Phải xác định đúng các lực tác dụng và chọn chiều dương trước khi lập phương trình.
\item (Sai) Có thể bỏ qua điều kiện áp dụng và chiều của lực mà kết quả vẫn luôn đúng.
\end{itemchoice}
}
\end{ex}

% ===== Câu 148 =====
% ID: AIQ_L10C3_FULL_0634
% Mức: VD
\begin{ex}
% ID: AIQ_L10C3_FULL_0634
% Mức: VD
Hai vật tiếp xúc cùng chuyển động, hợp lực ngoài làm hệ có gia tốc $1\,\text{m}$ /s^2. Tổng khối lượng hệ là $10\,\text{kg}$. Tính hợp lực ngoài tác dụng lên hệ.
\choiceTF
{\True Giá trị cần tìm là $10\,\text{N}$.}
{Giá trị cần tìm là $11\,\text{N}$.}
{\True Phải xác định đúng các lực tác dụng và chọn chiều dương trước khi lập phương trình.}
{Có thể bỏ qua điều kiện áp dụng và chiều của lực mà kết quả vẫn luôn đúng.}
\loigiai{
\begin{itemchoice}
\item (Đúng) Giá trị cần tìm là $10\,\text{N}$. (Vì): Xét toàn hệ, các lực tương tác trong triệt tiêu theo từng cặp; F=(m1+m2)a=10·1= $10\,\text{N}$. b) Sai vì sai giá trị. c) Đúng. d) Sai vì phải kiểm tra điều kiện áp dụng và chiều lực
\item (Sai) Giá trị cần tìm là $11\,\text{N}$.
\item (Đúng) Phải xác định đúng các lực tác dụng và chọn chiều dương trước khi lập phương trình.
\item (Sai) Có thể bỏ qua điều kiện áp dụng và chiều của lực mà kết quả vẫn luôn đúng.
\end{itemchoice}
}
\end{ex}

% ===== Câu 149 =====
% ID: AIQ_L10C3_FULL_0635
% Mức: VD
\begin{ex}
% ID: AIQ_L10C3_FULL_0635
% Mức: VD
Hai vật tiếp xúc cùng chuyển động, hợp lực ngoài làm hệ có gia tốc $2\,\text{m}$ /s^2. Tổng khối lượng hệ là $7\,\text{kg}$. Tính hợp lực ngoài tác dụng lên hệ.
\choiceTF
{\True Giá trị cần tìm là $14\,\text{N}$.}
{Giá trị cần tìm là $15\,\text{N}$.}
{\True Phải xác định đúng các lực tác dụng và chọn chiều dương trước khi lập phương trình.}
{Có thể bỏ qua điều kiện áp dụng và chiều của lực mà kết quả vẫn luôn đúng.}
\loigiai{
\begin{itemchoice}
\item (Đúng) Giá trị cần tìm là $14\,\text{N}$. (Vì): Xét toàn hệ, các lực tương tác trong triệt tiêu theo từng cặp; F=(m1+m2)a=7·2= $14\,\text{N}$. b) Sai vì sai giá trị. c) Đúng. d) Sai vì phải kiểm tra điều kiện áp dụng và chiều lực
\item (Sai) Giá trị cần tìm là $15\,\text{N}$.
\item (Đúng) Phải xác định đúng các lực tác dụng và chọn chiều dương trước khi lập phương trình.
\item (Sai) Có thể bỏ qua điều kiện áp dụng và chiều của lực mà kết quả vẫn luôn đúng.
\end{itemchoice}
}
\end{ex}

% ===== Câu 150 =====
% ID: AIQ_L10C3_FULL_0636
% Mức: VD
\begin{ex}
% ID: AIQ_L10C3_FULL_0636
% Mức: VD
Hai vật tiếp xúc cùng chuyển động, hợp lực ngoài làm hệ có gia tốc $3\,\text{m}$ /s^2. Tổng khối lượng hệ là $10\,\text{kg}$. Tính hợp lực ngoài tác dụng lên hệ.
\choiceTF
{\True Giá trị cần tìm là $30\,\text{N}$.}
{Giá trị cần tìm là $31\,\text{N}$.}
{\True Phải xác định đúng các lực tác dụng và chọn chiều dương trước khi lập phương trình.}
{Có thể bỏ qua điều kiện áp dụng và chiều của lực mà kết quả vẫn luôn đúng.}
\loigiai{
\begin{itemchoice}
\item (Đúng) Giá trị cần tìm là $30\,\text{N}$. (Vì): Xét toàn hệ, các lực tương tác trong triệt tiêu theo từng cặp; F=(m1+m2)a=10·3= $30\,\text{N}$. b) Sai vì sai giá trị. c) Đúng. d) Sai vì phải kiểm tra điều kiện áp dụng và chiều lực
\item (Sai) Giá trị cần tìm là $31\,\text{N}$.
\item (Đúng) Phải xác định đúng các lực tác dụng và chọn chiều dương trước khi lập phương trình.
\item (Sai) Có thể bỏ qua điều kiện áp dụng và chiều của lực mà kết quả vẫn luôn đúng.
\end{itemchoice}
}
\end{ex}

% ===== Câu 151 =====
% ID: AIQ_L10C3_FULL_0637
% Mức: TH
\begin{ex}
% ID: AIQ_L10C3_FULL_0637
% Mức: TH
Hai vật tiếp xúc cùng chuyển động, hợp lực ngoài làm hệ có gia tốc $4\,\text{m}$ /s^2. Tổng khối lượng hệ là $9\,\text{kg}$. Tính hợp lực ngoài tác dụng lên hệ.
\shortans{36}
\loigiai{
Xét toàn hệ, các lực tương tác trong triệt tiêu theo từng cặp; F=(m1+m2)a=9·4= $36\,\text{N}$.
}
\end{ex}

% ===== Câu 152 =====
% ID: AIQ_L10C3_FULL_0638
% Mức: TH
\begin{ex}
% ID: AIQ_L10C3_FULL_0638
% Mức: TH
Hai vật tiếp xúc cùng chuyển động, hợp lực ngoài làm hệ có gia tốc $1\,\text{m}$ /s^2. Tổng khối lượng hệ là $12\,\text{kg}$. Tính hợp lực ngoài tác dụng lên hệ.
\shortans{12}
\loigiai{
Xét toàn hệ, các lực tương tác trong triệt tiêu theo từng cặp; F=(m1+m2)a=12·1= $12\,\text{N}$.
}
\end{ex}

% ===== Câu 153 =====
% ID: AIQ_L10C3_FULL_0639
% Mức: VD
\begin{ex}
% ID: AIQ_L10C3_FULL_0639
% Mức: VD
Hai vật tiếp xúc cùng chuyển động, hợp lực ngoài làm hệ có gia tốc $2\,\text{m}$ /s^2. Tổng khối lượng hệ là $9\,\text{kg}$. Tính hợp lực ngoài tác dụng lên hệ.
\shortans{18}
\loigiai{
Xét toàn hệ, các lực tương tác trong triệt tiêu theo từng cặp; F=(m1+m2)a=9·2= $18\,\text{N}$.
}
\end{ex}

% ===== Câu 154 =====
% ID: AIQ_L10C3_FULL_0640
% Mức: VD
\begin{ex}
% ID: AIQ_L10C3_FULL_0640
% Mức: VD
Hai vật tiếp xúc cùng chuyển động, hợp lực ngoài làm hệ có gia tốc $3\,\text{m}$ /s^2. Tổng khối lượng hệ là $12\,\text{kg}$. Tính hợp lực ngoài tác dụng lên hệ.
\shortans{36}
\loigiai{
Xét toàn hệ, các lực tương tác trong triệt tiêu theo từng cặp; F=(m1+m2)a=12·3= $36\,\text{N}$.
}
\end{ex}

% ===== Câu 155 =====
% ID: AIQ_L10C3_FULL_0641
% Mức: VD
\begin{ex}
% ID: AIQ_L10C3_FULL_0641
% Mức: VD
Hai vật tiếp xúc cùng chuyển động, hợp lực ngoài làm hệ có gia tốc $4\,\text{m}$ /s^2. Tổng khối lượng hệ là $10\,\text{kg}$. Tính hợp lực ngoài tác dụng lên hệ.
\shortans{40}
\loigiai{
Xét toàn hệ, các lực tương tác trong triệt tiêu theo từng cặp; F=(m1+m2)a=10·4= $40\,\text{N}$.
}
\end{ex}

% ===== Câu 156 =====
% ID: AIQ_L10C3_FULL_0642
% Mức: VDC
\begin{ex}
% ID: AIQ_L10C3_FULL_0642
% Mức: VDC
Hai vật tiếp xúc cùng chuyển động, hợp lực ngoài làm hệ có gia tốc $1\,\text{m}$ /s^2. Tổng khối lượng hệ là $7\,\text{kg}$. Tính hợp lực ngoài tác dụng lên hệ.
\shortans{7}
\loigiai{
Xét toàn hệ, các lực tương tác trong triệt tiêu theo từng cặp; F=(m1+m2)a=7·1= $7\,\text{N}$.
}
\end{ex}

% ===== Câu 157 =====
% ID: AIQ_L10C3_FULL_0643
% Mức: VD
\begin{bt}
% ID: AIQ_L10C3_FULL_0643
% Mức: VD
Hai vật tiếp xúc cùng chuyển động, hợp lực ngoài làm hệ có gia tốc $1\,\text{m}$ /s^2. Tổng khối lượng hệ là $14\,\text{kg}$. Tính hợp lực ngoài tác dụng lên hệ.
\loigiai{
Phân tích lực, chọn hệ quy chiếu và áp dụng định luật hoặc quy tắc phù hợp. Xét toàn hệ, các lực tương tác trong triệt tiêu theo từng cặp; F=(m1+m2)a=14·1= $14\,\text{N}$.
}
\end{bt}

% ===== Câu 158 =====
% ID: AIQ_L10C3_FULL_0644
% Mức: VD
\begin{bt}
% ID: AIQ_L10C3_FULL_0644
% Mức: VD
Hai vật tiếp xúc cùng chuyển động, hợp lực ngoài làm hệ có gia tốc $2\,\text{m}$ /s^2. Tổng khối lượng hệ là $6\,\text{kg}$. Tính hợp lực ngoài tác dụng lên hệ.
\loigiai{
Phân tích lực, chọn hệ quy chiếu và áp dụng định luật hoặc quy tắc phù hợp. Xét toàn hệ, các lực tương tác trong triệt tiêu theo từng cặp; F=(m1+m2)a=6·2= $12\,\text{N}$.
}
\end{bt}

% ===== Câu 159 =====
% ID: AIQ_L10C3_FULL_0645
% Mức: VDC
\begin{bt}
% ID: AIQ_L10C3_FULL_0645
% Mức: VDC
Hai vật tiếp xúc cùng chuyển động, hợp lực ngoài làm hệ có gia tốc $3\,\text{m}$ /s^2. Tổng khối lượng hệ là $9\,\text{kg}$. Tính hợp lực ngoài tác dụng lên hệ.
\loigiai{
Phân tích lực, chọn hệ quy chiếu và áp dụng định luật hoặc quy tắc phù hợp. Xét toàn hệ, các lực tương tác trong triệt tiêu theo từng cặp; F=(m1+m2)a=9·3= $27\,\text{N}$.
}
\end{bt}

% ===== Câu 160 =====
% ID: AIQ_L10C3_FULL_0646
% Mức: VDC
\begin{bt}
% ID: AIQ_L10C3_FULL_0646
% Mức: VDC
Hai vật tiếp xúc cùng chuyển động, hợp lực ngoài làm hệ có gia tốc $4\,\text{m}$ /s^2. Tổng khối lượng hệ là $12\,\text{kg}$. Tính hợp lực ngoài tác dụng lên hệ.
\loigiai{
Phân tích lực, chọn hệ quy chiếu và áp dụng định luật hoặc quy tắc phù hợp. Xét toàn hệ, các lực tương tác trong triệt tiêu theo từng cặp; F=(m1+m2)a=12·4= $48\,\text{N}$.
}
\end{bt}

% ===== Câu 161 =====
% ID: AIQ_L10C3_FULL_0647
% Mức: VDC
\begin{bt}
% ID: AIQ_L10C3_FULL_0647
% Mức: VDC
Hai vật tiếp xúc cùng chuyển động, hợp lực ngoài làm hệ có gia tốc $1\,\text{m}$ /s^2. Tổng khối lượng hệ là $9\,\text{kg}$. Tính hợp lực ngoài tác dụng lên hệ.
\loigiai{
Phân tích lực, chọn hệ quy chiếu và áp dụng định luật hoặc quy tắc phù hợp. Xét toàn hệ, các lực tương tác trong triệt tiêu theo từng cặp; F=(m1+m2)a=9·1= $9\,\text{N}$.
}
\end{bt}

% ===== Câu 162 =====
% ID: AIQ_L10C3_FULL_0648
% Mức: VDC
\begin{bt}
% ID: AIQ_L10C3_FULL_0648
% Mức: VDC
Hai vật tiếp xúc cùng chuyển động, hợp lực ngoài làm hệ có gia tốc $2\,\text{m}$ /s^2. Tổng khối lượng hệ là $12\,\text{kg}$. Tính hợp lực ngoài tác dụng lên hệ.
\loigiai{
Phân tích lực, chọn hệ quy chiếu và áp dụng định luật hoặc quy tắc phù hợp. Xét toàn hệ, các lực tương tác trong triệt tiêu theo từng cặp; F=(m1+m2)a=12·2= $24\,\text{N}$.
}
\end{bt}
